\documentclass[10pt]{article}
\usepackage{amsmath, amssymb, amsthm}
\usepackage{graphicx}
\usepackage{booktabs}
\usepackage[utf8]{inputenc}
\usepackage{hyperref}
\usepackage{geometry}
\geometry{margin=1in}
\usepackage{float}
\usepackage{placeins}

\newtheorem{theorem}{Theorem}
\newtheorem{lemma}{Lemma}
\newtheorem{definition}{Definition}
\newtheorem{corollary}{Corollary}

\title{The Quantum Structural Triad: \\ Fluctuations, Entropy, and Correlations \\ as Interdependent Primitives}
\author{Ovidiu Tataru}
\date{}

\begin{document}

\maketitle
\begin{abstract}
We identify a primitive structural condition in quantum theory: energy fluctuations, von Neumann entropy, and quantum mutual information constitute an inseparable triad that cannot vanish on average in any physically realizable, locally coupled quantum system. This triad expresses the minimal structural integrity required for a quantum state to persist, evolve, and mediate phenomena. From locality and dynamical non-stationarity we derive two complementary inequalities—a Hamiltonian-capped and a speed-based form—that establish explicit quantitative thresholds for systemic viability. These thresholds reveal how fluctuations, entropy, and correlations must coexist and compensate one another to sustain quantum activity. The results apply to pure and mixed states, extend naturally to open systems via entropy production, and generalize to algebraic QFT. Beyond theory, the framework provides operational criteria for assessing structural viability in quantum technologies, from qubit certification to memory stability and error-correction design. While subsystem quantities are used for analytical resolution, the theory is fundamentally systemic: it diagnoses the coherence and persistence of the total quantum state rather than the autonomy of its parts. We further conjecture that analogous structural constraints may extend to globally topological quantum systems, whose persistent entropic and correlational features may reflect a generalized form of systemic viability.
\end{abstract}



\section{Introduction}

Quantum theory imposes fundamental constraints on what qualifies as a physically viable subsystem.
From the uncertainty principle to modern information-theoretic bounds\cite{NielsenChuang2010,CsiszarKorner2011},
it has long been clear that fluctuations, entropy, and correlations cannot all vanish simultaneously
without undermining the internal consistency of the theory. Each captures a distinct aspect of quantum structure:
energy fluctuations quantify sensitivity to local generators\cite{BraunsteinCaves1994,Escher2011},
von Neumann entropy measures local information content\cite{Blanco2013,Horodecki2009},
and mutual information encodes connectivity across subsystems\cite{Safranek2017,Frowis2018}.
While each has been widely studied in isolation or in pairs—
as resources for metrology\cite{Taddei2013,DeffnerCampbell2017},
thermodynamics\cite{BrandaoHorodecki2015,Spohn1978},
or information processing\cite{BreuerPetruccione2002}—
their collective interdependence has remained uncharacterized.

Existing work has established important pairwise constraints:
generalized speed limits connect quantum Fisher information to entropy production\cite{DeffnerLutz2013,NishiyamaHasegawa2024},
fluctuation--speed relations bound entropy growth in open systems\cite{Hamazaki2024,ShiraishiSaito2018},
and entanglement persists under noise through correlation–entropy trade-offs\cite{Benedetti2012,BenattiFloreanini2005}.
Yet none of these results capture the essential \emph{three-way structural interdependence}
that we show to be indispensable for subsystem viability.

Our central claim is that fluctuation, entropy, and correlation form a \emph{quantum structural triad}:
in any locally coupled, physically realizable \textbf{quantum system},
their product cannot vanish when averaged over time and across partitions.
Pointwise vanishing may occur, but systemic persistence is mandatory.
This defines a minimal structural threshold below which a subsystem becomes
\emph{quantum-structurally inert}\cite{DeffnerCampbell2017,NishiyamaHasegawa2025}:
incapable of sustaining dynamics, storing information, or mediating interactions.
In contrast, persistent nonvanishing of the triad guarantees that a subsystem
possesses the minimal structural integrity required to remain quantum-mechanically relevant.

The novelty lies not in recognizing that these three quantities matter,
but in deriving explicit, unit-correct \emph{quantitative thresholds}
that determine how much of each is required and how they trade off against one another.
These thresholds scale nontrivially with system size and minimal energy scales\cite{MargolusLevitin1998,Kliesch2014},
providing a falsifiable criterion for subsystem viability.
Unlike heuristic measures of coherence, our results emerge from first principles,
yielding both theoretical universality and practical tools for quantum technologies.


\paragraph{Problem statement.}
Consider a system on a bounded-degree finite-range interaction graph and any bipartition $(A|\bar A)$ crossed by at least one interaction. Define the triad $\Pi_A(t)=\sqrt{F_Q(\rho_A;H_A)}\,S_A\,I(A{:}\bar A)$. Without presupposing positivity of any leg, establish necessary non-inertness:
prove that the triad cannot vanish in time-average across the cut-set under locality, genuine cross-cut coupling, and global non-stationarity. \cite{NachtergaeleSims2010,NielsenChuang2010,BraunsteinCaves1994}
\section{Quantum Structural Triad: Definition and Role}

We define the \emph{quantum structural triad} of a subsystem $A$ in a global state $\rho(t)$ as
\begin{equation}
\Pi_A(t) = \sqrt{F_Q(\rho_A(t); H_A)} \; S_A(t) \; I(A:\bar{A},t),
\label{eq:triad}
\end{equation}
where:
\begin{itemize}
    \item $F_Q(\rho_A; H_A)$ is the quantum Fisher information of the reduced state $\rho_A$ with
    respect to the local generator $H_A$, equivalent to the Bures metric curvature at time $t$,
    capturing \emph{fluctuation strength}.
    \item $S_A(t) = -\mathrm{Tr}[\rho_A(t)\log \rho_A(t)]$ is the von Neumann entropy of the subsystem,
    quantifying \emph{local informational content}.
    \item $I(A:\bar{A},t) = S_A(t)+S_{\bar{A}}(t)-S_{A\bar{A}}(t)$ is the quantum mutual information,
    measuring total (quantum and classical) \emph{correlation} between $A$ and its complement.
\end{itemize}

Each leg reflects a distinct facet of quantum structure:
\begin{enumerate}
    \item $\sqrt{F_Q}$ encodes local dynamical sensitivity to perturbations,
    \item $S_A$ encodes intrinsic informational richness,
    \item $I(A:\bar{A})$ encodes connectivity to the rest of the system.
\end{enumerate}

\paragraph{Interpretation.}
We interpret $\Pi_A(t)$ as the \emph{instantaneous structural integrity} of subsystem $A$.
If $\Pi_A(t) = 0$, then $A$ is simultaneously fluctuation-free, pure, or decorrelated, and thus fails to
contribute meaningfully to the system's quantum structure.

\paragraph{Quantum-Structurally Inert (QSI).}
\begin{definition}[D2 (QSI on an epoch)]
A subsystem $A$ is \emph{quantum-structurally inert} on a measurable epoch $K\subset[0,\infty)$ with $\delta(K)>0$ if $\Pi_A(t)=0$ for almost every $t\in K$.
\end{definition}

\paragraph{Systemic Triad.}
We define the systemic measure as the sum over all subsystems:
\begin{equation}
\Pi(t) = \sum_{A} \Pi_A(t).
\end{equation}
Our main result is that $\Pi(t)$, while possibly vanishing at isolated instants, cannot vanish on average
across subsystems and over time in any locally coupled, physically viable quantum system. This
establishes the triad as a \emph{minimal structural threshold for subsystem viability}.
\section{Hypotheses and Objects}

\paragraph{Conventions.}
All logarithms are natural; entropies are measured in nats. For any subsystem $A$, let $d_A$ denote the
Hilbert space dimension of $H_A$ and $d_{\bar{A}}$ that of its complement.

\paragraph{Interaction graph and cut-set.}
We represent the system by a bounded-degree, finite-range interaction graph $G=(V,E)$ associated
with the local terms of the Hamiltonian. A cut-set $E$ consists of the nearest-neighbor edges crossing
a bipartition $(A|\bar{A})$. Its cardinality is denoted $|E|$, and for intensive averages we use
\begin{equation}
\frac{1}{|E|} \sum_{A \in E} (\cdots).
\end{equation}
\paragraph{Structural Hypotheses (H1–H3).}
To state our results precisely, we impose three physically natural hypotheses:
\begin{description}
    \item[H1. Locality.] The Hamiltonian is a sum of local terms with finite interaction range and uniformly bounded degree on the interaction graph $G$.
    \item[H2. Genuine coupling.] Each bipartition defined by the cut-set $E$ includes at least one nonzero coupling term, preventing complete isolation of subsystems.
    \item[H3. Global non-stationarity.] There exists a measurable set $K \subset [0,\infty)$ with lower density $\delta > 0$ such that $\dot\rho(t) \neq 0$ for almost every $t \in K$ (closed systems), or such that the CPTP evolution $\Lambda_t$ is non-stationary on $K$ (open systems).
\end{description}

\paragraph{Continuity caveat.}
Rank-change times can accumulate in engineered cases. Our results apply on \emph{full-rank epochs}, where $\sqrt{F_Q}$, $S_A$, and $I(A:\bar{A})$ are continuous (see Lemma~\ref{lemma:continuity}).

\paragraph{Dimension bounds.}
For any bipartition $(A,\bar{A})$, the mutual information is bounded by:
\begin{equation}
I(A:\bar{A}) \leq 2 \log \min(d_A, d_{\bar{A}}).
\end{equation}

\paragraph{Interpretation.}
These hypotheses exclude trivial or degenerate regimes (e.g., completely isolated subsystems, frozen dynamics, or rank-changing discontinuities). Within this admissible domain, we establish that the triad $\Pi_A$ defines a nonvanishing \emph{structural threshold} for subsystem viability.

\paragraph{Scope Note.}
Although expressed through subsystem observables, the framework is not subsystem-reductionist. The triad thresholds are intended to characterize the \emph{global} viability of the quantum state, with partitions introduced as analytical tools for structural resolution. We make no ontological claim about subsystems as independently real; they serve as a lens through which system-level structure becomes operationally visible.

\section{Geometric Uncertainty: QSL from Bures Angle}

The Mandelstam--Tamm form of the quantum speed limit provides a direct link between distinguishability
in state space and local energy fluctuations. For a subsystem $A$, the minimal time $\tau_A$ required
to evolve by a Bures angle $L$ obeys
\begin{equation}
\tau_A \geq \frac{\hbar L}{\Delta E_A},
\end{equation}
where $\Delta E_A$ is the time-averaged standard deviation of the local Hamiltonian $H_A$ in the
state $\rho_A$.

Equivalently,
\begin{equation}
\Delta E_A \geq \frac{\hbar L}{\tau_A}.
\end{equation}

\paragraph{Interpretation.}
This expresses the \emph{geometric face of uncertainty}: whenever the subsystem undergoes a
``distinguishable change'' (nonzero Bures angle), local energy fluctuations must be strictly positive.
In structural terms, no subsystem can remain indefinitely fluctuation-free while still evolving in
Hilbert space. This provides the first leg of the structural triad.
\section{Entropic Response: First Law of Entanglement}

For small perturbations of a quantum state, the first law of entanglement entropy provides a direct link
between entropy variation and modular Hamiltonian expectation values:
\begin{equation}
\delta S_A = \delta \langle H^{\mathrm{mod}}_A \rangle,
\end{equation}
where $H^{\mathrm{mod}}_A$ is the modular Hamiltonian associated with region $A$.

\paragraph{Interpretation.}
This relation shows that entropy cannot remain frozen under nontrivial local energy exchange.
Whenever the subsystem participates in dynamics, changes in $S_A$ are unavoidable. Entropy thus forms
the second leg of the structural triad, encoding the local informational footprint of the subsystem.

\paragraph{Caveat.}
The first law of entanglement is exact for small variations around a reference state, and is most
transparent when modular Hamiltonians are explicitly known (e.g., conformal field theory or holographic
settings). For general mixed states, it should be interpreted as a local differential relation rather than a
global identity. Nevertheless, it captures the essential fact: local entropy cannot vanish persistently in
a coupled, evolving quantum system.
\section{Correlation as Intrinsic}

The quantum mutual information across a bipartition $(A|\bar{A})$ is defined as
\begin{equation}
I(A:\bar{A}) = S_A + S_{\bar{A}} - S_{A\bar{A}}.
\end{equation}

It measures the total correlation—both classical and quantum—between the subsystem $A$ and its
complement $\bar{A}$.

\paragraph{Interpretation.}
Entropy is never a purely local quantity: it necessarily reflects correlations with the environment or
complementary subsystems. Thus mutual information, as the third leg of the triad, is structurally
unavoidable in any nontrivial quantum regime. If correlations vanished identically across a genuine cut,
the subsystem would be decoupled and structurally irrelevant.

\paragraph{Structural role.}
Correlation ensures that local fluctuations and entropy are not isolated artifacts but are embedded
in a larger relational structure. Together with $\sqrt{F_Q}$ and $S_A$, mutual information completes
the minimal triad of structural viability.
\section{Per-Subsystem Inequality}

Combining the fluctuation, entropy, and correlation terms, and using the instantaneous Bures
speed bound, we obtain on any active interval:
\begin{equation}
\Delta E_A \, S_A \, I(A:\bar{A}) \;\;\geq\;\; \frac{\hbar}{2} \, \dot{L}_B \, S_A \, I(A:\bar{A}),
\end{equation}
where $\dot{L}_B$ is the instantaneous Bures angle rate.

Integrating over time and applying the density-to-time-average lemma on a measurable active set $K$
with lower density $\delta > 0$, we arrive at
\begin{equation}
\langle \Pi_A \rangle_t \;\;\geq\;\; \delta \, L \, \frac{S_A I(A:\bar{A})}{\tau_A}.
\end{equation}

\paragraph{Interpretation.}
This inequality expresses the \emph{structural inseparability} of the triad. Variance, entropy, and
correlation cannot collapse simultaneously for a genuinely coupled subsystem. At least one leg must
sustain a strictly positive floor over active epochs, ensuring that $\Pi_A$ cannot vanish in time average.

\paragraph{Structural consequence.}
Whenever $\langle \Pi_A \rangle_t = 0$, the subsystem is quantum-structurally inert: it contributes no
persistent fluctuations, entropy, or correlations. The inequality therefore defines a \emph{per-subsystem
structural threshold for viability}.
\section{Systemic Inequality}

Summing over all subsystems yields the global triad inequality:
\begin{equation}
\Pi(t) \;=\; \sum_{A} \sqrt{F_Q(\rho_A(t); H_A)} \, S_A(t) \, I(A:\bar{A},t)
\;\;\geq\;\; \sum_{A} \frac{L \, S_A I(A:\bar{A})}{\tau_A}.
\end{equation}

As long as the trajectory in state space has nonzero Bures length ($L>0$), the right-hand side is
strictly positive.

\paragraph{Interpretation.}
Whenever the global state undergoes distinguishable evolution, the triad enforces a \emph{system-wide
structural threshold}. Even if individual subsystems fluctuate between relevance and irrelevance at
specific instants, the ensemble cannot sustain a collectively vanishing product across time.

\paragraph{Structural consequence.}
This inequality establishes that no physically viable, locally coupled quantum system can remain
entirely fluctuationless, unentangled, and pure across all subsystems. Systemic persistence of the
triad is thus a universal requirement for maintaining quantum structural integrity.
\section{Intensive Formulation and Root Inequalities}

\subsection{Normalized Intensive Triad}

To compare across systems and isolate structure from scale, we define a dimensionless intensive version
of the triad for each cut $(A|\bar{A})$:
\begin{equation}
\widehat{\Psi}_A(t)
:= \underbrace{\frac{\hbar\,\sqrt{F_Q(\rho_A(t);H_A)}}{2\|H_A\|_{\mathrm{op}}}}_{\in[0,1]}
\;\cdot\;
\underbrace{\frac{S_A(t)}{\log d_A}}_{\in[0,1]}
\;\cdot\;
\underbrace{\frac{I(A:\bar{A},t)}{2\log\min(d_A,d_{\bar{A}})}}_{\in[0,1]}.
\label{eq:intensive-triad}
\end{equation}
Each factor is bounded by standard inequalities, so $0\le \widehat{\Psi}_A(t)\le 1$.
We use the cut-set average and its time average:
\[
\widehat{\Psi}_{\mathrm{cut}}(t) := \frac{1}{|E|}\sum_{A\in E}\widehat{\Psi}_A(t),
\qquad
\left\langle \widehat{\Psi}_{\mathrm{cut}} \right\rangle_t
:= \liminf_{T\to\infty}\frac{1}{T}\int_0^T \widehat{\Psi}_{\mathrm{cut}}(t)\,dt.
\]

\subsection{Hamiltonian-Capped Structural Threshold}

\begin{theorem}[Hamiltonian-Capped Structural Threshold]
\label{thm:ham-cap}
Under H1–H2 (no speed assumptions), for any cut-set $E$,
\begin{equation}
\frac{1}{|E|}\sum_{A\in E} \Delta E_A\,S_A\,I(A:\bar{A})
\;\;\ge\;\;
2\,(\log d_{\min})^2\,\|H\|_{\min}\;
\Big\langle \widehat{\Psi}_{\mathrm{cut}} \Big\rangle_t,
\qquad
d_{\min}:=\min_{A\in E}\min(d_A,d_{\bar{A}}),
\label{eq:ham-cap-struct-threshold}
\end{equation}
where $\|H\|_{\min}:=\min_{A\in E}\|H_A\|_{\mathrm{op}}$.
\end{theorem}

\paragraph{Proof sketch.}
From \eqref{eq:intensive-triad},
\[
\sqrt{F_Q}\,S_A\,I(A:\bar{A})
:= \frac{2\|H_A\|_{\mathrm{op}}}{\hbar}\,(\log d_A)\,[2\log\min(d_A,d_{\bar{A}})]\;\widehat{\Psi}_A(t).
\]
Use $\Delta E_A\ge \frac{\hbar}{2}\sqrt{F_Q}$ to replace $\sqrt{F_Q}$, then average over $A\in E$ and
lower-bound $\|H_A\|_{\mathrm{op}}, \log d_A, \log\min(d_A,d_{\bar{A}})$ by their respective minima.
Time-averaging yields \eqref{eq:ham-cap-struct-threshold}. \qed

\paragraph{Interpretation.}
This threshold is \emph{unconditional in speed}: even without assuming a dynamical floor, the normalized
structural integrity $\langle \widehat{\Psi}_{\mathrm{cut}}\rangle_t$ forces a positive lower bound set by the
\emph{minimal local energy scale} and \emph{minimal subsystem dimension}. It quantifies how structural
integrity trades off against hardware scales.

\subsection{Speed-Based Structural Threshold}

\begin{theorem}[Speed-Based Structural Threshold]
\label{thm:speed-based}
Assume H1–H2–H3. By Lemmas L2–L4, there exist $\varepsilon_F,\varepsilon_S,\varepsilon_I>0$ and a measurable subset $K^\star\subset K$ with $\delta(K^\star)>0$ such that
\[
\sqrt{F_Q(\rho_A(t);H_A)}\ge \varepsilon_F,\quad
S_A(t)\ge \varepsilon_S,\quad
I(A:\bar{A},t)\ge \varepsilon_I
\]
for almost every $t\in K^\star$ and all $A\in E$. Then
\begin{equation}
\frac{1}{|E|}\sum_{A\in E} \Delta E_A\,S_A\,I(A:\bar{A})
\;\;\ge\;\; \frac{\hbar}{2}\;\delta(K^\star)\;\varepsilon_F\,\varepsilon_S\,\varepsilon_I.
\label{eq:speed-based-struct-threshold}
\end{equation}
This uses the density-to-time-average lemma.
\end{theorem}

\paragraph{Interpretation.}
This threshold is \emph{dynamical}: it converts a \emph{Bures/QFI speed floor} on a derived active epoch $K^\star\subset K$ into a
strictly positive structural lower bound, with explicit constants $(\hbar/2)\,\delta(K^\star)\,\varepsilon_F\,\varepsilon_S\,\varepsilon_I$.
It is typically tighter during nonequilibrium evolution. Finite intersections of positive-measure sets retain positive measure; density-to-time-average then yields the stated bound.

\subsection{Trade-offs and Scaling}
Both thresholds quantify \emph{trade-offs}: if one leg (fluctuation, entropy, or correlation) is small, the
others must compensate to keep the left-hand side above threshold. System-size enters via $(\log d_{\min})^2$:
larger minimal Hilbert spaces demand stronger combined resources to maintain structural integrity.
\section{Technical Lemmas}

The following lemmas establish the auxiliary bounds needed to derive the structural thresholds.

\begin{lemma}[QFI–variance bound]
For any density operator $\rho$ and Hermitian generator $H$,
\begin{equation}
F_Q(\rho;H) \;\leq\; \frac{4}{\hbar^2}\,\mathrm{Var}_\rho(H),
\qquad
\Delta E \;\geq\; \frac{\hbar}{2}\sqrt{F_Q(\rho;H)}.
\end{equation}
\label{lemma:qfi-var}
\end{lemma}

\begin{lemma}[Instantaneous Bures speed]
For unitary dynamics generated by $H_A$,
\begin{equation}
\dot{L}_B(t) = \tfrac{1}{2}\sqrt{F_Q(\rho_A(t); H_A)}.
\end{equation}
For general CPTP dynamics, $\dot{L}_B(t)\le \tfrac{1}{2}\sqrt{F_Q}$.
\label{lemma:bures-speed}
\end{lemma}

\begin{lemma}[Continuity and full-rank epochs]
For finite-dimensional piecewise $C^1$ trajectories, the functions
$t\mapsto S_A(t)$, $t\mapsto I(A:\bar{A},t)$, and
$t\mapsto \sqrt{F_Q(\rho_A(t);H_A)}$
are continuous on intervals where $\rho_A(t)$ is full rank.
Rank-change times form a set of measure zero.
\label{lemma:continuity}
\end{lemma}

\begin{definition}[D1 (Lower density)]
For measurable $K\subseteq [0,\infty)$, define
\begin{equation}
\delta(K) := \liminf_{T\to\infty} \frac{|K\cap [0,T]|}{T}.
\end{equation}
\end{definition}

\begin{lemma}[Density-to-time-average]
Let $f:[0,\infty)\to [0,\infty)$ be measurable, and suppose $f(t)\ge \alpha>0$
for all $t\in K$ with $\delta(K)>0$. Then
\begin{equation}
\liminf_{T\to\infty}\frac{1}{T}\int_0^T f(t)\,dt \;\;\geq\;\; \delta(K)\,\alpha.
\end{equation}
\label{lemma:density}
\end{lemma}

\paragraph{Structural Lemmas (L1–L4).}
\begin{lemma}[L1: Reduced motion from global motion]
Under H1–H2–H3, reduced dynamics is non-stationary for almost every time in $K$ except possibly a null set: fine-tuned decoupling across a genuine cut has measure zero.
\end{lemma}
\begin{lemma}[L2: Fluctuation leg]
From non-stationarity on $K$ and Lemma~\ref{lemma:bures-speed}, there exists a subset $K_F\subset K$ with $\delta(K_F)>0$ and $\sqrt{F_Q(\rho_A;H_A)}>0$ for almost every $t\in K_F$.
\end{lemma}
\begin{lemma}[L3: Correlation leg]
With a cross-cut coupling (H2) and locality (H1), Lieb–Robinson propagation implies $I(A:\bar A)>0$ after a microscopic time on a positive-measure subset $K_I\subset K$.
\end{lemma}
\begin{lemma}[L4: Entropy leg]
For closed pure global states, $I=2S_A$. For mixed/open dynamics, Spohn's inequality implies strictly positive local entropy production away from steady state on a positive-measure subset $K_S\subset K$.
\end{lemma}

\section{Discussion: Structural Relevance and Irrelevance}

\subsection{Complementary Thresholds}
The two inequalities derived above provide complementary perspectives on subsystem viability:
\begin{itemize}
    \item \textbf{Hamiltonian-capped threshold:} unconditional, tied to the smallest operator norm
    $\|H\|_{\min}$ and subsystem dimension $d_{\min}$. It reflects the \emph{structural potential}
    inherent in the system's composition.
    \item \textbf{Speed-based threshold:} conditional on a nonzero Bures/QFI speed floor, it converts
    dynamical activity into a strictly positive structural lower bound, typically tighter in
    nonequilibrium regimes.
\end{itemize}
These represent the only two physically consistent, unit-correct routes to nonvanishing lower bounds on
the fluctuation–entropy–correlation product.

\subsection{Units and Physical Meaning}
The structural triad combines distinct units:
\begin{itemize}
    \item $\sqrt{F_Q}$ carries units of inverse time (distinguishability rate),
    \item $\Delta E$ carries energy,
    \item $S$ and $I$ are dimensionless (nats).
\end{itemize}
The Hamiltonian-capped threshold introduces an energy scale via $\|H\|_{\min}$,
while the speed-based threshold introduces a temporal scale via the Bures speed.
In both cases, the inequality enforces that the product of these resources
cannot vanish under admissible conditions.

\subsection{Boundary Cases and Structural Irrelevance}
The thresholds become trivial precisely when an \emph{excluded hypothesis} holds:
\begin{enumerate}
    \item $I\equiv 0$ (violates H2: no genuine cross-cut coupling),
    \item $\sqrt{F_Q}\equiv 0$ or zero-speed protocols (violates H3: global non-stationarity),
    \item Rank-changing discontinuities beyond measure-zero sets (covered by continuity caveat).
\end{enumerate}
In each case, the subsystem is \emph{quantum-structurally inert} (QSI): incapable of acting as a
reliable mediator, memory, or participant in nontrivial dynamics.

\subsection{Trade-offs and Design Rules}
The inequalities highlight explicit trade-offs:
\begin{itemize}
    \item Small entropy can be compensated by larger fluctuations or correlations,
    \item Weak correlations demand stronger fluctuations and entropy,
    \item Minimal subsystem dimension $d_{\min}$ sets a rising structural cost with scale.
\end{itemize}
Thus subsystem viability is governed not by any single resource but by their interdependence.
This suggests practical design rules: viable subsystems must maintain all three legs above
thresholds in some combination, and resource deficits in one leg require reinforcement in the others.

\subsection{Experimental Accessibility}
Few-body experiments can probe the triad directly via:
\begin{itemize}
    \item QFI estimates from Bures angle trajectories,
    \item Entropy via tomography or entropic witnesses,
    \item Mutual information via correlation measurements.
\end{itemize}
Many-body systems require witness-based or coarse estimators, but our finite-window QFI bound
(Appendix) enables positive certification even with partial data.

\subsection*{Corollary (no admissible QSI epochs)}
If H1–H2–H3 hold, then $A$ cannot be QSI on $K$: otherwise $\Pi_A=0$ almost everywhere on a positive-density epoch $K^\star$ contradicts the thresholds above.

\section{Quantum Computing Applications}

The structural thresholds provide practical, physics-based exclusion criteria for quantum technologies.
Because they are \emph{necessary conditions}, any subsystem falling below threshold is
\emph{guaranteed} to lack quantum functionality, allowing definitive elimination of nonviable qubits,
gates, or memories.

\subsection{Qubit Viability Certification}
For each physical qubit $i$, measure or estimate
\begin{equation}
\Pi_i = \sqrt{F_Q(\rho_i;H_i)} \, S_i \, I(i:\bar{i}).
\end{equation}
\noindent
A qubit is structurally viable only if
\[
\Pi_i \;\geq\; \max\{\text{Hamiltonian-capped threshold},\;\text{Speed-based threshold}\}.
\]
Qubits with $\Pi_i$ below threshold are quantum-structurally inert and cannot support reliable
quantum logic.

\subsection{Error Correction Prioritization}
Error correction is resource-intensive. Our framework identifies qubits worth correcting:
\[
\text{Apply QEC only if } \Pi_i \geq \text{threshold}.
\]
Qubits below threshold are structurally irrelevant and cannot be salvaged by error correction,
saving resources in large-scale devices.

\subsection{Quantum Memory Lifetime}
The speed-based inequality provides a bound on memory coherence:
\begin{equation}
\text{Quantum memory is viable while }
\delta(t)\,\varepsilon_F(t)\,\varepsilon_S(t)\,\varepsilon_I(t) \;\geq\; \frac{2}{\hbar}\,\Pi_{\min}.
\end{equation}
When any component falls below its floor, the memory is structurally incapable of preserving
quantum coherence.

\subsection{Gate Operation Certification}
For multi-qubit gates, structural viability requires
\begin{equation}
\Pi_{\text{gate}} \;\geq\; \frac{\hbar}{2}\,\delta_{\text{gate}}\,
\varepsilon_F \varepsilon_S \varepsilon_I
\end{equation}
throughout the gate execution. Gates failing this criterion cannot realize genuine quantum logic,
regardless of classical fidelity measures.

\subsection{NISQ Device Optimization}
In noisy intermediate-scale devices, the thresholds enable:
\begin{itemize}
    \item identification of qubits that are ``quantum enough'' to include in algorithms,
    \item bounding of maximum circuit depth before triad decay below threshold,
    \item optimization of connectivity patterns based on correlation persistence.
\end{itemize}

\subsection{Towards Principled Design}
Current practice in quantum engineering is heuristic, guided by coherence times and error rates.
Our framework introduces principled design rules grounded in first-principles physics.
It answers questions such as:
\begin{itemize}
    \item ``Is this qubit structurally capable of quantum computation?''
    \item ``When should error correction be abandoned?''
    \item ``What are the fundamental limits of this processor architecture?''
\end{itemize}
By providing definitive exclusion criteria, the structural triad transforms quantum engineering
from empirical tuning to threshold-driven design.

\section{Open-System Extension: Structural Relevance Under Dissipation}

\subsection{Setting and Hypotheses}
We extend the structural triad to dissipative dynamics governed by a local Lindbladian
\[
\mathcal{L} = \sum_X \mathcal{L}_X,
\]
with finite interaction range and bounded degree. The CPTP propagator is
$\Lambda_t = e^{t\mathcal{L}}$ (or time-periodic for driven steady states).
Cuts are defined by the same nearest-neighbor cut-set $E$.
Each cut $(A|\bar{A})$ is assumed genuinely coupled by at least one Lindblad or Hamiltonian term.
Dissipative Lieb–Robinson bounds ensure finite-speed propagation of influence.

\paragraph{H3 (open).}
There exists a measurable epoch $K\subset[0,\infty)$ of non-stationary CPTP evolution (away from a fixed point) with $\delta(K)>0$.

\subsection{Open-System Triad Reformulation}
Even under dissipation, the triadic structure persists in generalized form:
\begin{itemize}
    \item \textbf{Fluctuations:} The QFI leg is replaced by a channel speed measure
    $F_{\mathrm{chan}}$, defined via monotone Riemannian metrics on CPTP maps.
    \item \textbf{Entropy:} The entropy leg becomes the local entropy production rate
    \[
    \dot{S}_{\mathrm{prod},A}(t) := -\frac{d}{dt} D(\rho_A(t)\|\sigma_A) \geq 0,
    \]
    where $\sigma_A$ is the full-rank local steady state (Spohn's inequality), strictly $>0$ away from steady state.
    \item \textbf{Correlations:} The correlation leg is extended to system–environment
    mutual information, accessible via Stinespring dilations.
\end{itemize}

\subsection{Intensive Open-System Triad}
Define the normalized intensive open-system triad:
\begin{equation}
\widehat{\Psi}^{\mathrm{open}}_A(t)
:= \underbrace{\frac{\sqrt{F_{\mathrm{chan}}(\Lambda_t|A)}}{\sqrt{F_{\max}}}}_{\in[0,1]}
\cdot
\underbrace{\frac{\dot{S}_{\mathrm{prod},A}(t)}{\dot{S}_{\max}}}_{\in[0,1]}
\cdot
\underbrace{\frac{I(\mathrm{sys}:\mathrm{rest},t)}{2\log d_{\min}}}_{\in[0,1]}.
\end{equation}
The cut-set intensive average is
\[
\widehat{\Psi}^{\mathrm{open}}_{\mathrm{cut}}(t) = \frac{1}{|E|}\sum_{A\in E}\widehat{\Psi}^{\mathrm{open}}_A(t),
\qquad
\overline{\widehat{\Psi}^{\mathrm{open}}_{\mathrm{cut}}}
:= \liminf_{T\to\infty}\frac{1}{T}\int_0^T \widehat{\Psi}^{\mathrm{open}}_{\mathrm{cut}}(t)\,dt.
\]

\begin{theorem}[Open-System Structural Threshold]
\label{thm:open}
Under locality, genuine coupling, finite-speed dissipative propagation, and existence of a non-stationary epoch $K$ with density $\delta>0$ and leg-wise positivity away from steady state,
\begin{equation}
\overline{\widehat{\Psi}^{\mathrm{open}}_{\mathrm{cut}}}
\;\;\geq\;\;
\delta \,
\frac{\varepsilon_F}{\sqrt{F_{\max}}}
\cdot
\frac{\varepsilon_S}{\dot{S}_{\max}}
\cdot
\frac{\varepsilon_I}{2\log d_{\min}}
\;>\;0.
\end{equation}
\end{theorem}

\paragraph{Interpretation.}
Even with dissipation, subsystems cannot persist in a state of structural irrelevance
except at true fixed points. Any nontrivial open-system evolution enforces persistent
nonvanishing of fluctuation, entropy production, and correlations, maintaining
quantum structural integrity in the presence of noise and dissipation.


\section{Conclusions}

We have identified a primitive structural condition in quantum theory: the inseparable triad of fluctuations, entropy, and correlations defines the minimal requirement for subsystem viability. No physically realizable, locally coupled quantum system can maintain dynamical relevance while any leg of the triad vanishes on average. The triad thus constitutes the fundamental expression of quantum structural integrity—the capacity to persist, evolve, and mediate phenomena.

From locality and dynamical non-stationarity we derived two complementary quantitative thresholds:
\begin{itemize}
    \item a \textbf{Hamiltonian-capped} inequality, governed by minimal local energy scales and subsystem dimensions,
    \item a \textbf{speed-based} inequality, governed by Bures evolution rates and entropic/correlation floors.
\end{itemize}
Together, these bounds reveal that fluctuations, entropy, and correlations are structurally interdependent. Any persistent deficit in one component must be balanced by the others, while the simultaneous vanishing of all three marks structural inertness. This inseparability defines the threshold between inert mathematical possibility and physically participatory reality.

Although the framework is expressed through subsystem observables—reduced fluctuations, local entropy, and mutual information—it is not reductionist. Subsystems function only as analytical constructs introduced to resolve structure within the total system. The triad theory is fundamentally systemic: it characterizes the integrity of the full quantum state, employing partitions as diagnostic tools rather than ontological divisions.

Subsystems falling below the derived thresholds are structurally inert: incapable of storing information, mediating interactions, or sustaining distinguishable evolution. In contrast, persistent triadic activity guarantees the minimal conditions for quantum persistence. This exclusion principle applies to pure and mixed states, extends naturally to Lindbladian open systems through entropy production, and generalizes to algebraic QFT.

Beyond foundational implications, the triad provides first-principles benchmarks for quantum technologies—criteria for qubit viability, memory stability, gate certification, and error-correction prioritization. It redefines quantum engineering in terms of structural necessity rather than empirical performance.

At a deeper level, the persistence of triadic activity suggests an informational origin of spacetime geometry. In modular-Hamiltonian and holographic formulations, the same structural relations reproduce the linearized Einstein equations and encode effective gravitational couplings. Gravity thereby emerges not as an additional force, but as the geometric manifestation of sustained quantum distinguishability and correlation—the curvature of persistent informational structure.

Future work will extend the computational annexes, formalize experimental witness protocols, and explore the link between subsystem structure, modular dynamics, and emergent geometry. Through this lens, quantum viability is not an operational convention but a universal principle of physical structure.

\paragraph{Conjectural Extension.}
While the present results rest on locality and dynamical evolution, we conjecture that analogous triadic thresholds may govern globally topological quantum systems, such as the toric code on a torus. Even without conventional dynamics, such systems exhibit persistent entropy, entanglement, and fluctuation under logical observables, suggesting that a generalized form of the structural triad may still signal systemic viability beyond local frameworks.
\section{Conjecture: Topological Structural Triad (Toric Code Case)}
\label{sec:topo-triad}

The structural triad, as derived above, presupposes locality and dynamical non-stationarity (H1–H3). 
However, certain quantum systems---notably the toric code---exhibit structural persistence in the 
absence of both. This motivates a conjectural extension of the framework to globally topological, 
non-dynamical quantum orders.

\subsection*{Mathematical Structure}

The toric code Hamiltonian on an $L \times L$ torus,
\begin{equation}
H = -J_e \sum_v A_v - J_m \sum_p B_p,
\end{equation}
acts on $N = 2L^2$ qubits located on the edges of the lattice. 
Here $A_v = \prod_{i \in v} \sigma^x_i$ and $B_p = \prod_{i \in p} \sigma^z_i$ 
are the mutually commuting star and plaquette stabilizers. 
Despite its static ground-space structure, the toric code retains finite entropy, 
mutual information, and logical fluctuations under nonlocal observables.

\paragraph{Ground-space structure.}
All stabilizers commute and satisfy $A_v^2 = B_p^2 = I$, yielding 
\begin{equation}
\dim \mathcal{H}_{\mathrm{GS}} = 4,
\end{equation}
with logical operators $W_x, W_z$ forming noncontractible Wilson loops that anticommute on the torus.

\paragraph{Logical fluctuations.}
For the mixed ground-state density operator,
\begin{equation}
\rho_{\mathrm{mix}} = \frac{1}{4} \sum_{i=1}^{4} |\psi_i\rangle \langle \psi_i|,
\end{equation}
the logical Wilson loops satisfy
\begin{equation}
\langle W_x \rangle = \langle W_z \rangle = 0, 
\qquad 
\Delta W_x = \Delta W_z = 1,
\end{equation}
indicating persistent logical uncertainty even in the absence of dynamical evolution.

\paragraph{Entropic and correlational structure.}
For half-torus bipartitions at $L=2$,
\begin{align}
S_A &\approx 2.77~\text{nats}, \\
I(A:\bar{A}) &\approx 4.16,
\end{align}
leading to the observed topological triad values
\begin{equation}
\Pi^{\mathrm{topo}}_{x,z} 
  = \Delta W_{x,z} \, S_A \, I(A:\bar{A})
  \approx 11.53.
\end{equation}
The equality $\Pi^{\mathrm{topo}}_x = \Pi^{\mathrm{topo}}_z$ reflects isotropy under $J_e=J_m$.

\subsection*{Interpretation}

Even without dynamical evolution or local coupling, the triad remains finite. 
The logical fluctuations $\Delta W_{x,z}$ substitute for local energy variance, while 
$S_A$ and $I(A:\bar{A})$ preserve their informational roles. 
The toric code therefore satisfies a generalized static inequality,
\begin{equation}
\Pi^{\mathrm{topo}} 
   = \Delta_{\mathrm{logical}} \, S_A \, I(A:\bar{A}) > 0,
\end{equation}
expressing the structural viability of topological quantum systems.

\subsection*{Implications}

\begin{itemize}
    \item The persistence of $\Pi^{\mathrm{topo}}>0$ demonstrates that 
          systemic viability can be maintained through global logical structure, 
          not only through dynamical coupling.
    \item Logical operators act as ``virtual dynamics,'' sustaining distinguishability 
          and informational curvature in a stationary code space.
    \item This extension suggests that the triad characterizes 
          \emph{systemic structural persistence}---the capacity of a quantum system 
          to remain informationally active---whether sustained by local motion 
          or by topological coherence.
\end{itemize}

\subsection*{Future Computations}

To formalize this conjecture, subsequent work will:
\begin{itemize}
    \item Extend simulations to $L=3,4$ using stabilizer algebra rather than dense diagonalization;
    \item Evaluate the scaling of $S_A$, $I(A:\bar{A})$, and $\Pi^{\mathrm{topo}}$ with system size $L$;
    \item Test whether $\Pi^{\mathrm{topo}}/L^2$ approaches a constant, defining a topological invariant 
          of structural integrity.
\end{itemize}

This conjectural extension positions the triad as a universal descriptor of quantum structural 
persistence, valid across both local and topological phases.
\section{Open Discussion: Qualitative Outlook}

The structural triad developed in this work—linking fluctuations, entropy, and correlations—
suggests a deeper informational origin of spacetime geometry.
The persistence of nonzero triadic activity implies that no physically realizable subsystem can remain
simultaneously fluctuationless, pure, and decorrelated.
When examined within conformal field theories admitting an Einstein gravity dual,
this persistence principle yields a correspondence between
\emph{quantum structural integrity} and \emph{geometric curvature}.

This correspondence follows from modular–Hamiltonian thermodynamics and information geometry.
The modular first law of Casini, Huerta, and Myers\cite{CasiniHuertaMyers2011}
relates variations of entanglement entropy to local energy density,
while the second variation of relative entropy defines the Bures
(symmetric–logarithmic–derivative) quantum Fisher metric
\cite{BraunsteinCaves1994,Petz1996,Hayashi2006},
which quantifies curvature in the space of quantum states.
Together these relations imply that informational and physical curvature obey parallel structural constraints,
so persistent triadic activity guarantees a nondegenerate informational metric and, by extension,
a nontrivial spacetime geometry.

In this view, gravitational coupling and cosmological scaling emerge from the triad’s intrinsic dynamics
rather than from fundamental inputs.
Gravity thus appears not as an additional force but as a manifestation of persistent quantum
distinguishability and correlation.
A detailed mathematical development of this correspondence—including the derivation of effective
Einstein–Hilbert dynamics and the renormalization-group calibration of the gravitational constants—
is presented in the companion work\cite{FaulknerLewkowyczMaldacena2013,Jacobson1995,BrownHenneaux1986}.
\begin{center}
\fbox{%
\begin{minipage}{0.93\linewidth}
\small
\textbf{Illustrative Correspondence: From Quantum Structure to Curvature.} \\
\medskip
Formally, the second variation of relative entropy,
\[
\delta^{2}S(\rho\|\sigma)=\tfrac{1}{4}F_{Q},
\]
and the linearized Einstein constraint,
\[
\delta^{2}S_{B}\propto
\int_{B}\!\xi^{\mu}
\big(E_{\mu\nu}-8\pi G\,T_{\mu\nu}\big)\xi^{\nu},
\]
share the same structural pattern: a curvature term balanced by an
information–energy term. This parallel—first recognized in the
modular–Hamiltonian approach of Casini, Huerta, and Myers, and later
developed by Faulkner, Lewkowycz, and Maldacena—shows that the geometry
of state space and the geometry of spacetime obey homologous variational
laws.

\medskip
Dimensionally, an effective coupling
\[
G^{-1}\sim\frac{\hbar}{c^{5}}
(\text{fluctuation rate})^{2}
(\text{entropy})
(\text{correlation}),
\]
follows the same scaling as the triadic threshold derived earlier.
In this view, spacetime curvature measures the
\emph{persistent distinguishability} of local quantum configurations:
a perfectly inert field (vanishing triad) corresponds to a flat
informational geometry.
\end{minipage}}
\end{center}


\section*{Acknowledgements} 

This project began with a long-standing question: how does the principle \textit{natura abhorret vacuum}—the idea that physical reality never admits absolute emptiness—retain its validity from the microscopic to the cosmological scale?
To explore that question I carried out a symbolic reduction of representative equations drawn from quantum, thermodynamic, relativistic, and cosmological theories.
When all non-universal or model-specific quantities were removed, three elements remained invariant across those domains: energy fluctuations, entropy, and Planck’s constant.
In pure states, entropy manifests as entanglement entropy, revealing the same structure in its quantum form.
These invariants became the foundation of the structural triad developed and verified in this work.

Large-language-model systems were employed as analytical and computational assistants. To limit both human and algorithmic bias and to control potential hallucination, all theoretical results were subsequently verified with the \emph{Rocq} proof assistant and cross-validated through independent numerical simulations. 

The full verification materials, including Rocq proof scripts, computational notebooks, and simulation outputs, are available from the author’s GitHub repository upon request. 

Although not a professional physicist or formally affiliated researcher, I carried out the conceptual development, interpretation, and synthesis independently. The involvement of AI and proof-assistant tools should be understood as a form of assisted reasoning, not as authorship. I retain complete epistemic responsibility for the arguments, methods, and conclusions advanced here and accept responsibility for any remaining errors.

\appendix
% Annex transition and overview
\section*{Annexes: Computational and Formal Backing}
The annexes compile notation and normalization details, auxiliary lemmas and proofs, and the full
reproducibility package: computational validations A1--A13 alongside the Rocq/Coq verification.
For navigation, see the annex index at Section~\ref{app:annex}.
\addcontentsline{toc}{section}{Annexes: Computational and Formal Backing}

% This annex is included after \appendix in the main file.
% It provides technical proofs, normalization checks, scaling relations,
% and a concise computational validation summary backing the main results.

% This annex is included after \appendix in the main file.
% It provides technical proofs, normalization checks, scaling relations,
% and a concise computational validation summary backing the main results.

% Unnumbered theorem environment used below
\theoremstyle{plain}
\newtheorem*{theorem*}{Theorem}

% --- Begin inlined: Annex_Notation_Units ---
\section{Notation, Units, and Normalizations}

We recall the subsystem triad (Eq.~\eqref{eq:triad})
\[
\Pi_A(t) = \sqrt{F_Q(\rho_A(t);H_A)}\, S_A(t)\, I(A{:}\bar{A},t),
\]
and its intensive normalized form (Eq.~\eqref{eq:intensive-triad})
\[
\widehat{\Psi}_A(t)
:::= \frac{\hbar\,\sqrt{F_Q(\rho_A;H_A)}}{2\lVert H_A\rVert_{\mathrm{op}}}
\cdot \frac{S_A}{\log d_A}
\cdot \frac{I(A{:}\bar{A})}{2\log\min(d_A,d_{\bar{A}})}\in[0,1].
\]

\paragraph{Units.}
With time encoding, \(\sqrt{F_Q}\) has units of inverse time, \(\Delta E\) has units of energy, while \(S\) and \(I\) are dimensionless (nats).
The Hamiltonian-capped inequality introduces an energy scale via \(\lVert H_A\rVert_{\mathrm{op}}\), and the speed-based form introduces a temporal scale via Bures/QFI speed.

\paragraph{Normalization bounds.}
(i) \(\hbar\sqrt{F_Q}/(2\lVert H_A\rVert_{\mathrm{op}})\le 1\) (unitary time encoding; cf. Lemma~\ref{lemma:qfi-var}).
(ii) \(S_A\le \log d_A\).
(iii) \(I(A{:}\bar{A})\le 2\log\min(d_A,d_{\bar{A}})\).
Hence \(\widehat{\Psi}_A\in[0,1]\).
% --- End inlined: Annex_Notation_Units ---

% --- Begin inlined: Annex_Lemmas ---
\section{Proof Details for Auxiliary Lemmas}

\begin{lemma}[QFI–variance bound; Lemma~\ref{lemma:qfi-var}]
For any density operator \(\rho\) and Hermitian generator \(H\),
\(F_Q(\rho;H) \le 4\,\mathrm{Var}_\rho(H)/\hbar^2\) and thus \(\Delta E\ge (\hbar/2)\sqrt{F_Q}\).
\end{lemma}
\begin{proof}
Using the SLD representation \(F_Q(\rho;H)=\mathrm{Tr}(\rho L^2)\) where \(\partial_\theta\rho = (L\rho+\rho L)/2\) for the unitary encoding \(\rho(\theta)=e^{-\tfrac{i}{\hbar}\theta H}\rho\,e^{\tfrac{i}{\hbar}\theta H}\), one has
\(\partial_\theta\rho = -\tfrac{i}{\hbar}[H,\rho]\). Monotonicity of the Bures metric and Cauchy–Schwarz yield the standard inequality \(F_Q\le 4\,\mathrm{Var}_\rho(H)/\hbar^2\) (see \cite{BraunsteinCaves1994,Hayashi2006}). Taking square roots gives \(\Delta E\ge (\hbar/2)\sqrt{F_Q}\).
\end{proof}

\begin{lemma}[Instantaneous Bures speed; Lemma~\ref{lemma:bures-speed}]
For unitary dynamics on \(A\), \(\dot L_B(t)=\tfrac{1}{2}\sqrt{F_Q(\rho_A(t);H_A)}\). For general CPTP dynamics, \(\dot L_B\le \tfrac{1}{2}\sqrt{F_Q}\).
\end{lemma}
\begin{proof}
The Bures angle rate equals half the SLD metric speed \cite{BraunsteinCaves1994}. For unitary encoding this is an equality; for CPTP maps, contractivity of monotone metrics implies an upper bound.
\end{proof}

\begin{lemma}[Continuity and full-rank epochs; Lemma~\ref{lemma:continuity}]
On intervals where \(\rho_A(t)\) is full rank and piecewise \(C^1\), the maps \(t\mapsto S_A(t)\), \(t\mapsto I(A{:}\bar A,t)\), and \(t\mapsto \sqrt{F_Q}\) are continuous; rank-change times have measure zero.
\end{lemma}
\begin{proof}
Spectral continuity of finite-dimensional density matrices implies continuity of von Neumann entropy away from eigenvalue collisions; mutual information is a linear combination of entropies. Standard perturbation theory and continuity of the SLD metric on full-rank states yield continuity of \(\sqrt{F_Q}\). Rank changes occur at isolated times in piecewise-smooth trajectories, forming a null set.
\end{proof}

\begin{lemma}[Density-to-time-average; Lemma~\ref{lemma:density}]
If \(f\ge 0\) and \(f(t)\ge \alpha>0\) for all \(t\in K\) with lower density \(\delta(K)>0\), then \(\liminf_{T\to\infty}\tfrac{1}{T}\int_0^T f\,dt\ge \delta(K)\,\alpha\).
\end{lemma}
\begin{proof}
For any \(T>0\), \(\int_0^T f\,dt\ge \int_{K\cap[0,T]} \! f\,dt\ge \alpha\,|K\cap[0,T]|\). Divide by \(T\) and take \(\liminf\) to obtain the claim.
\end{proof}

\paragraph{Structural Lemmas L1–L4.}
Under locality (H1) and genuine cross-cut coupling (H2), Lieb–Robinson bounds imply finite-speed propagation and, after a microscopic time, strictly positive mutual information across a genuine cut on a positive-measure subset (L3) \cite{NachtergaeleSims2010}. Global non-stationarity (H3) excludes protocols with vanishing Bures speed on a positive-density epoch, ensuring a positive-speed subset for \(\sqrt{F_Q}\) (L2). For closed pure states, \(I=2S_A\) implies an entropy floor whenever correlations are present; for open systems, Spohn’s inequality yields strictly positive local entropy production away from steady state (L4) \cite{Spohn1978}. Reduced non-stationarity (L1) follows generically from global motion except possibly on a null set.

\section{Proofs of the Main Inequalities}

\begin{theorem*}[Hamiltonian-Capped Structural Threshold; Thm.~\ref{thm:ham-cap}]
For any cut-set \(E\) under H1–H2,
\[
\frac{1}{|E|}\sum_{A\in E} \Delta E_A\,S_A\,I(A{:}\bar A)
\;\ge\; 2(\log d_{\min})^2\,\lVert H\rVert_{\min}\;\big\langle \widehat{\Psi}_{\mathrm{cut}}\big\rangle_t.
\]
\end{theorem*}
\begin{proof}
From Eq.~\eqref{eq:intensive-triad},
\(\sqrt{F_Q}\,S_A\,I = \tfrac{2\lVert H_A\rVert_{\mathrm{op}}}{\hbar}\,(\log d_A)\,[2\log\min(d_A,d_{\bar A})] \,\widehat{\Psi}_A\).
By Lemma~\ref{lemma:qfi-var}, \(\Delta E_A\ge (\hbar/2)\sqrt{F_Q}\). Hence
\[
\Delta E_A\,S_A\,I \;\ge\; \lVert H_A\rVert_{\mathrm{op}}\,(\log d_A)\,[2\log\min(d_A,d_{\bar A})] \,\widehat{\Psi}_A.
\]
Averaging over \(A\in E\) and lower-bounding by \(\lVert H\rVert_{\min}\), \(\log d_{\min}\) yields the claim after time averaging.
\end{proof}

\begin{theorem*}[Speed-Based Structural Threshold; Thm.~\ref{thm:speed-based}]
Under H1–H2–H3, there exist \(\varepsilon_F,\varepsilon_S,\varepsilon_I>0\) and \(K^*\subset K\) with \(\delta(K^*)>0\) such that
\[
\frac{1}{|E|}\sum_{A\in E} \Delta E_A\,S_A\,I(A{:}\bar A)\;\ge\; \frac{\hbar}{2}\,\delta(K^*)\,\varepsilon_F\,\varepsilon_S\,\varepsilon_I.
\]
\end{theorem*}
\begin{proof}
On \(K^*\), the leg-wise floors hold almost everywhere. By Lemma~\ref{lemma:qfi-var}, pointwise \(\Delta E_A\,S_A\,I\ge (\hbar/2)\,\sqrt{F_Q}\,S_A\,I\ge (\hbar/2)\,\varepsilon_F\,\varepsilon_S\,\varepsilon_I\). Apply Lemma~\ref{lemma:density} to the nonnegative function \(f(t)=(1/|E|)\sum_A \Delta E_A S_A I\) with \(K=K^*\).
\end{proof}

\section{Cut-Set Averaging and Scaling}
Let \(E\) be the set of nearest-neighbor edges crossing a bipartition. The intensive average
\(\widehat{\Psi}_{\mathrm{cut}}(t)=\tfrac{1}{|E|}\sum_{A\in E}\widehat{\Psi}_A(t)\) is dimensionless and \(\in[0,1]\).
The Hamiltonian-capped inequality scales with \((\log d_{\min})^2\,\lVert H\rVert_{\min}\), making explicit how minimal subsystem dimension and minimal local energy scale set the structural cost of viability. For lattice systems with fixed on-site dimension, the factor grows with the size of the minimal balanced cut, matching the observed growth of raw \(\Pi\) and the stabilization of intensive measures recorded in the simulations.

\section{Open-System Extension: Details}
For local Lindbladians \(\mathcal L=\sum_X\mathcal L_X\) with finite range and bounded degree, define a channel speed leg \(F_{\mathrm{chan}}\) via a contractive monotone Riemannian metric \cite{Petz1996}. Let
\[
\widehat{\Psi}^{\mathrm{open}}_A
::= \frac{\sqrt{F_{\mathrm{chan}}(\Lambda_t|A)}}{\sqrt{F_{\max}}}
\cdot \frac{\dot S_{\mathrm{prod},A}}{\dot S_{\max}}\cdot\frac{I(\mathrm{sys}{:}\mathrm{rest})}{2\log d_{\min}}\in[0,1].
\]
Spohn’s inequality \cite{Spohn1978} yields \(\dot S_{\mathrm{prod},A}\ge 0\), strictly \(>0\) away from steady state. Dissipative Lieb–Robinson bounds provide finite-speed influence. With leg-wise positivity on a subset \(K\) of density \(\delta>0\), the time-average of the cut-set intensive measure obeys Theorem~\ref{thm:open} as stated in the main text.


\section{Formal Verification (Rocq/Coq)}

The derivations underlying the structural thresholds have been mechanized in the Rocq/Coq development under \texttt{Hyppopotamus/rocq-project}. The formalization mirrors the paper’s structure (notation, normalization, product lifts) and includes the modular triad architecture (system interface, leg floors, normalization, sum lifting). Build and checking use the generated Coq makefiles.

\paragraph{Build.}
From the repository root (build inside the clean tree):
\begin{verbatim}
make -C Hyppopotamus/rocq-project/src_clean -f CoqMakefile -j"$(nproc)" | cat
\end{verbatim}

\paragraph{Logical checking (coqchk).}
Core check (recommended):
\begin{verbatim}
# Legacy constructive + triad variant
coqchk -R Hyppopotamus/rocq-project/src_clean ESSEClean \
  ESSEClean.ESSEBoxed ESSEClean.ESSEList ESSEClean.PerBoundClean \
  ESSEClean.PinskerClean ESSEClean.PinskerSquaredClean ESSEClean.GapFloorClean \
  ESSEClean.AveragingPerA ESSEClean.PiClean ESSEClean.Positivity \
  ESSEClean.PurityClean ESSEClean.FisherClean ESSEClean.EntropyClean \
  ESSEClean.TimeAvgClean ESSEClean.ESSEUniversal ESSEClean.ESSEUniversal_structure

# Optional: modular triad theory (excluding HCap while WIP)
coqchk -R Hyppopotamus/rocq-project/src_clean ESSEClean \
  ESSEClean.theory.SystemSpec ESSEClean.theory.TriadSignals \
  ESSEClean.theory.Bridges ESSEClean.theory.LegFloors \
  ESSEClean.theory.Normalization ESSEClean.theory.SumLifting \
  ESSEClean.theory.SpeedBound \
  ESSEClean.open.Channels ESSEClean.open.SpeedOpen \
  ESSEClean.tests.TriadExamples
\end{verbatim}
Successful runs return exit code 0 and produce no output beyond warnings. The mapping from paper sections to modules is summarized in the project \texttt{README.md} and \texttt{REPRODUCIBILITY.md}. The Hamiltonian-capped modular lemma (\texttt{theory/HCap.v}) is under active development and intentionally omitted from the optional check here.

\paragraph{Environment and Artifact.}
For long-term reproducibility, we recommend pinning a toolchain and providing a frozen artifact:
\begin{itemize}
  \item Toolchain (example via opam):
\begin{verbatim}
opam switch create esse-coq ocaml-4.14.1
opam install coq coq-mathcomp-ssreflect coq-coquelicot coq-elpi
# then run the build and coqchk commands above
\end{verbatim}
  \item Container (optional): ship a Dockerfile that runs \verb|make| and \verb|coqchk| on the pinned versions.
  \item Frozen release: publish a GitHub release (or DOI archive) containing source, _CoqProject, generated CoqMakefile, and precompiled \verb|.vo| files so \verb|coqchk| can be executed without rebuilding.
  \item REPL usage: 
\begin{verbatim}
coqtop -R Hyppopotamus/rocq-project/src_clean ESSEClean
(* inside Coq *)
Require Import ESSEClean.theory.SystemSpec ESSEClean.theory.Normalization.
\end{verbatim}
\end{itemize}

\paragraph{Formal Coverage (under stated hypotheses).}
The formal development proves the core triad theory under the manuscript's hypotheses (H1--H3, continuity, finite-dimensional cuts):
\begin{itemize}
  \item \textbf{System primitives}: \texttt{theory/SystemSpec.v} (admissible systems, minima, and log-monotonicity; e.g., \(\ln d_{\min}>0\), \(|E|>0\)).
  \item \textbf{Normalization}: \texttt{theory/Normalization.v} (\(0\le\widehat{\Psi}_A\le 1\) given positivity of scales).
  \item \textbf{Sum/monotone lifts}: \texttt{theory/SumLifting.v} (triple product monotonicity and sum monotonicity used throughout).
  \item \textbf{Triad signals (state-agnostic)}: \texttt{ESSEUniversal\_structure.v} (per-cut triad \(\Psi_A\), good-epoch floors, averaging lemmas, witness-based floors). This file contains the fully proved, state-agnostic core corresponding to the paper’s discrete-time triad arguments.
  \item \textbf{Bridges/leg floors stubs}: \texttt{theory/Bridges.v}, \texttt{theory/LegFloors.v} (interfaces used to connect physical witnesses to leg-wise floors).
  \item \textbf{Speed route}: \texttt{theory/SpeedBound.v} currently provides a compiling placeholder; a strengthened lemma can replace it without impacting upstream structure.
  \item \textbf{Hamiltonian-capped route}: \texttt{theory/HCap.v} is work-in-progress; its omission does not affect the proved state-agnostic triad floors.
\end{itemize}
All results are checked with \verb|coqchk| per the commands above; no additional “limitations” are introduced beyond the paper’s hypotheses.

\section{Reproducibility Runbook (Scripts)}

All commands are intended to be run from the repository root. Outputs are written under \texttt{Hyppopotamus/Scripts Demonstration/plots} or \texttt{.../results}.

\begin{itemize}
  \item A1 (core engine):\newline
  \verb|python3 -u "Hyppopotamus/Scripts Demonstration/1_esse_core_simulation.py"|
  \item A2 (pure-state validation):\newline
  \verb|python3 -u "Hyppopotamus/Scripts Demonstration/2_esse_pure_state_validation.py"|
  \item A3 (comprehensive + partitions):\newline
  \verb|python3 -u "Hyppopotamus/Scripts Demonstration/3_esse_comprehensive_validation.py"|
  \item A4 (TEBD Néel + $\tau$):\newline
  \verb|python3 -u "Hyppopotamus/Scripts Demonstration/4_run_neel_tebd_tau.py"|
  \item A5 (TEBD per-$A$ metrics):\newline
  \verb|python3 -u "Hyppopotamus/Scripts Demonstration/5_run_tebd_perA_metrics.py"|
  \item A6 (delta ladder TEBD):\newline
  \verb|python3 -u "Hyppopotamus/Scripts Demonstration/6_run_delta_ladder_tebd.py"|
  \item A7 (commuting vs non-commuting):\newline
  \verb|python3 -u "Hyppopotamus/Scripts Demonstration/7_esse_commuting_vs_noncommuting.py"|
  \item A8 (ETH eigenstate validation):\newline
  \verb|python3 -u "Hyppopotamus/Scripts Demonstration/8_esse_eigenstate_validation.py"|
  \item A9 (ball-$K_A$ proxy):\newline
  \verb|python3 -u "Hyppopotamus/Scripts Demonstration/9_esse_ball_KA_proxy.py"|
  \item A10 (triad intersection density):\newline
  \verb|python3 -u "Hyppopotamus/Scripts Demonstration/10_triad_density_test.py"|
  \item A11 (mixed-state quench stress):\newline
  \verb|python3 -u "Hyppopotamus/Scripts Demonstration/11_esse_mixed_quench_stress.py"|
  \item A12 (weak interaction sweep):\newline
  \verb|python3 -u "Hyppopotamus/Scripts Demonstration/12_esse_weak_interaction_test.py"|
  \item A13 (integrable XX validation):\newline
  \verb|python3 -u "Hyppopotamus/Scripts Demonstration/13_esse_integrable_validation.py"|
    \item A14 (topological triad; toric code):\newline
  \verb|python3 -u "Hyppopotamus/Scripts Demonstration/14_toric_topo_triad.py" --L 2 --state mixed-gs|
\end{itemize}

These runs reproduce the figures and tables referenced in the standalone computation appendices and support the inequalities proved above. For large TEBD configurations, adjust \verb|--steps| and \verb|--chi| to fit available memory.

\paragraph{Annex curation note.}
Scripts emit raw artifacts under \texttt{Hyppopotamus/Scripts Demonstration/plots} and \texttt{.../results}. For the annex, we copy the curated outputs into \texttt{Hyppopotamus/Annex results/A\#} with standardized filenames used by the figure/table references above.

\section{Annex Results (Scripts A1--A13)}\label{app:annex}

This section collects figure placeholders and file references corresponding to the reproducible runs. Paths point to the repository locations under \texttt{Hyppopotamus/Annex results/}. Figures use file paths with spaces via \verb|\detokenize{...}|.

\paragraph{Reading guide.}
Each A\# subsection includes a core metrics table with columns: $\langle\Pi\rangle$ (or $\langle\Pi_{\mathrm{rate}}\rangle$ when rates are reported), $\tau_A^4$, RHS$_{\mathrm{cap}}$, RHS$_{\mathrm{speed}}$, and $\langle\widehat{\Psi}\rangle$. N/A indicates not applicable or not reported.
\\[2pt]\textit{Units:} $\langle\Pi\rangle$ and $\langle\Pi_{\mathrm{rate}}\rangle$ dimensionless (rate when noted); $\tau_A^4$ in time units of the simulation; RHS$_{\mathrm{cap}}$ and RHS$_{\mathrm{speed}}$ in energy-like units per context; $\langle\widehat{\Psi}\rangle$ dimensionless in $[0,1]$.
\\[2pt]\textit{Note:} $\tau_A^4$ denotes a fourth-moment time proxy used in TEBD diagnostics; theoretical bounds use $\tau_A$ (Bures-time) but we report $\tau_A^4$ where available as a robust summary.

\subsection*{Annex Index (A1--A13)}
\begin{itemize}
  \item A1: Core engine (non-commuting vs commuting), Section~\ref{app:a1}.
  \item A2: Pure-state validation (trial variability and bounds), Section~\ref{app:a2}.
  \item A3: Comprehensive validation and partition stress tests, Section~\ref{app:a3}.
  \item A4: TEBD N\'eel state with $\tau$ sampling, Section~\ref{app:a4}.
  \item A5: TEBD per-$A$ metrics, Section~\ref{app:a5}.
  \item A6: Delta ladder TEBD, Section~\ref{app:a6}.
  \item A7: Commuting vs non-commuting comparison, Section~\ref{app:a7}.
  \item A8: ETH eigenstate validation, Section~\ref{app:a8}.
  \item A9: Ball-$K_A$ proxy, Section~\ref{app:a9}.
  \item A10: Triad intersection density, Section~\ref{app:a10}.
  \item A11: Mixed-state quench stress, Section~\ref{app:a11}.
  \item A12: Weak interaction sweep, Section~\ref{app:a12}.
  \item A13: Integrable XX validation, Section~\ref{app:a13}.
\end{itemize}

\subsection{A1: Core engine (non-commuting vs commuting)}\label{app:a1}
\begin{figure}[H] % or [!ht] if you prefer floating
  \centering
  \includegraphics[width=0.75\linewidth]{A1_noncomm_N6.png}
  \caption{A1: Core engine diagnostics (non-commuting Hamiltonian).}
  \label{fig:A1_noncomm}
\end{figure}

\begin{figure}[H]
  \centering
  \includegraphics[width=0.75\linewidth]{A1_comm_N6.png}
  \caption{A1: Core engine diagnostics (commuting Hamiltonian).}
  \label{fig:A1_comm}
\end{figure}
\FloatBarrier
\noindent CSVs: \texttt{A1/A1\_summary.csv}. Log: \texttt{A1/A1\_run.log}.

\paragraph{Summary.}
Non-commuting: $\langle\Pi_{\mathrm{rate}}\rangle=1.157$, LHS($\sum_A\Delta E\,S\,I$)$=1.258$, RHS$_{\mathrm{cap}}=0.0787$, RHS$_{\mathrm{speed}}=0.0150$, $r\approx0.385$.\\
Commuting: $\langle\Pi_{\mathrm{rate}}\rangle=0.870$, LHS$=2.504$, RHS$_{\mathrm{cap}}=0.0570$, RHS$_{\mathrm{speed}}\approx7.4\times10^{-5}$, $r\approx0.933$.

\begin{table}[!ht]
\centering
\caption{A1: per-case summary}
\label{tab:a1-summary}
\begin{tabular}{@{}lccccc@{}}
\toprule
Case & $\langle\Pi_{\mathrm{rate}}\rangle$ & LHS & RHS$_{\mathrm{cap}}$ & RHS$_{\mathrm{speed}}$ & $r$ \\
\midrule
non-comm & 1.157 & 1.258 & 0.0787 & 0.0150 & 0.385 \\
comm     & 0.870 & 2.504 & 0.0570 & $7.4\!\times\!10^{-5}$ & 0.933 \\
\bottomrule
\end{tabular}
\end{table}

\begin{table}[!ht]
\centering
\caption{A1: core metrics}
\label{tab:a1-core}
\begin{tabular}{@{}lccccc@{}}
\toprule
Case & $\langle\Pi_{\mathrm{rate}}\rangle$ & $\tau_A^4$ & RHS$_{\mathrm{cap}}$ & RHS$_{\mathrm{speed}}$ & $\langle\widehat{\Psi}\rangle$ \\
\midrule
non-comm & 1.157 & — & 0.0787 & 0.0150 & 0.061 \\
comm     & 0.870 & — & 0.0570 & $7.4\times10^{-5}$ & 0.059 \\
\bottomrule
\end{tabular}
\end{table}

% (moved A2/A3/A8/A9 core tables into their corresponding sections)

\subsection{A2: Pure-state validation (trial variability and bounds)}\label{app:a2}


\paragraph{Summary.}
Two trials per case (product and Haar) at $N{=}6$; CSV reports per-trial $\langle\Pi_{\mathrm{rate}}\rangle$, liminf proxies, and thresholds; plots saved under \texttt{plots/pure/}.

\begin{table}[!ht]
\centering
\caption{A2: condensed head/tail (pure-state CSV)}
\label{tab:a2-condensed}
\begin{tabular}{@{}lccccc@{}}
\toprule
Case & $\langle\Pi_{\mathrm{rate}}\rangle$ & $\tau_A^4$ & RHS$_{\mathrm{cap}}$ & RHS$_{\mathrm{speed}}$ & $\langle\widehat{\Psi}\rangle$ \\
\midrule
noncomm & 1.1570 & N/A & 0.0787 & 0.0150 & N/A \\
comm    & 0.8703 & N/A & 0.0570 & $7.4\times10^{-5}$ & N/A \\
\bottomrule
\end{tabular}
\end{table}

\begin{table}[!ht]
\centering
\caption{A2: auxiliary metrics (from summary CSV)}
\label{tab:a2-aux}
\begin{tabular}{@{}lcccc@{}}
\toprule
 & $\varepsilon_F$ & $\varepsilon_S$ & $\varepsilon_I$ & RHS$_{\mathrm{speed}}$ \\
\midrule
aggregate & 4.689 & 2.810 & 5.620 & 6.170 \\
\bottomrule
\end{tabular}
\end{table}

% (moved A10 core table into its section)

\subsection{A3: Comprehensive validation and partition stress tests}\label{app:a3}
No figure placeholder needed. Artifact references: \texttt{A3/A3\_liminf\_1site.csv}, \texttt{A3/A3\_partition\_stress.csv}, \texttt{A3/A3\_thresholds\_summary.csv}. Extract notes: \texttt{A3/A3\_liminf\_extract.txt}. Log: \texttt{A3/A3\_run.log}.

\paragraph{Summary (from log).}
1-site non-commuting: final liminf mean $1.528\pm0.304$; thresholds: RHS$_{\mathrm{cap}}=1.074\times10^{-1}$, RHS$_{\mathrm{speed}}=1.162\times10^{-2}$, $\langle\widehat{\Psi}\rangle=0.0831$.\\
2-site non-commuting: final liminf mean $7.967\pm4.148$; thresholds: RHS$_{\mathrm{cap}}=1.265$, RHS$_{\mathrm{speed}}=0.368$, $\langle\widehat{\Psi}\rangle=0.121$.\\
3-site non-commuting: final liminf mean $19.926\pm11.958$; thresholds: RHS$_{\mathrm{cap}}=4.998$, RHS$_{\mathrm{speed}}=2.561$, $\langle\widehat{\Psi}\rangle=0.138$.

\begin{table}[!ht]
\centering
\caption{A3: partition stress (means)}
\label{tab:a3-partitions}
\begin{tabular}{@{}lccc@{}}
\toprule
Partition & RHS$_{\mathrm{cap}}$ & RHS$_{\mathrm{speed}}$ & $\langle\widehat{\Psi}\rangle$ \\
\midrule
1-site & $1.074\!\times\!10^{-1}$ & $1.162\!\times\!10^{-2}$ & 0.0831 \\
2-site & $1.265$ & $0.368$ & 0.121 \\
3-site & $4.998$ & $2.561$ & 0.138 \\
\bottomrule
\end{tabular}
\end{table}


\begin{table}[!ht]
\centering
\caption{A3: 1-site liminf (non-comm vs comm)}
\label{tab:a3-1site}
\begin{tabular}{@{}lccc@{}}
\toprule
Case & $\langle\Pi\rangle_{\mathrm{liminf}}$ (mean) & std & conv-time (mean) \\
\midrule
non-comm & 1.528 & 0.304 & 6.62 \\
comm     & 1.865 & 0.101 & 5.45 \\
\bottomrule
\end{tabular}
\end{table}

\subsection{A4: TEBD N\'eel state with \(\tau\) sampling}\label{app:a4}
\begin{figure}[H]
  \centering
  \includegraphics[width=0.75\linewidth]{esse_tdmrg_neel_tau_L8_ell3_chi128_dt0.05_steps50_k2_plot.png}
  \caption{A4: TEBD diagnostics at cut size $\ell{=}3$.}
  \label{fig:A4_ell3}
\end{figure}
\FloatBarrier
\noindent CSVs: \texttt{A4/A4\_summary.csv}. Logs: \texttt{A4/A4\_run\_small.log}, \texttt{A4/A4\_run\_large.log}.

\begin{figure}[H]
  \centering
  \includegraphics[width=0.75\linewidth]{L8_ell4_chi128_dt0.05_steps300_k2_plot.png}
  \caption{A4: TEBD diagnostics at cut size $\ell{=}4$.}
  \label{fig:A4_ell4}
\end{figure}
\FloatBarrier
\noindent CSVs: \texttt{A4/A4\_summary.csv}. Logs: \texttt{A4/A4\_run\_small.log}, \texttt{A4/A4\_run\_large.log}.

\paragraph{Summary.}
Per-cut TEBD on N\'eel initial states at \(\ell{=}3,4\); CSV summarizes $\langle\Pi\rangle$ and $\tau_A^4$; figures show representative time traces.

\begin{table}[!ht]
\centering
\caption{A4: condensed head/tail (N\,=\,8)}
\label{tab:a4-condensed}
\begin{tabular}{@{}lccccc@{}}
\toprule
$\ell$ & $\langle\Pi\rangle$ & $\tau_A^4$ & RHS$_{\mathrm{cap}}$ & RHS$_{\mathrm{speed}}$ & $\langle\widehat{\Psi}\rangle$ \\
\midrule
3 & 4.988 & 3.604 & 2.494 & 0.100 & 0.0560 \\
4 & 18.060 & 9.027 & 9.030 & 4.841 & 0.0868 \\
\bottomrule
\end{tabular}
\end{table}




\subsection{A5: TEBD per-\(A\) metrics}\label{app:a5}
Per-cut metrics in CSV/JSON: \texttt{A5/esse\_tdmrg\_perA\_L8\_ell3\_perA.csv}, \texttt{A5/esse\_tdmrg\_perA\_L12\_ell6\_perA.csv} (and corresponding \texttt{.json}). Summaries: \texttt{A5/A5\_perA\_run1.csv}, \texttt{A5/A5\_perA\_run2.csv}. Logs: \texttt{A5/A5\_run1.log}, \texttt{A5/A5\_run2.log}.

\begin{table}[!ht]
\centering
\caption{A5: condensed metrics (Run1/Run2)\footnotesize~(RHScap computed as $\langle\Pi\rangle/2$ for TEBD per-$A$ runs)}
\label{tab:a5-condensed}
\begin{tabular}{@{}lccccc@{}}
\toprule
 & $\langle\Pi\rangle$ & $\tau_A^4$ & RHS$_{\mathrm{cap}}$ & RHS$_{\mathrm{speed}}$ & $\varepsilon_F$ \\
\midrule
Run1 & 7.554 & N/A & 3.777 & 1.284 & 1.303 \\
Run2 & 49.729 & N/A & 24.864 & 6.170 & 4.689 \\
\bottomrule
\end{tabular}
\end{table}

\begin{table}[!ht]
\centering
\caption{A5: Run2 liminf (non-comm vs comm)}
\label{tab:a5-run2-liminf}
\begin{tabular}{@{}lccc@{}}
\toprule
Case & $\langle\Pi\rangle_{\mathrm{liminf}}$ (mean) & std & conv-time (mean) \\
\midrule
non-comm & 1.528 & 0.304 & 6.62 \\
comm     & 1.865 & 0.101 & 5.45 \\
\bottomrule
\end{tabular}
\end{table}

\begin{table}[!ht]
\centering
\caption{A5: per-A head/tail (examples)}
\label{tab:a5-perA}
\begin{tabular}{@{}lccc@{}}
\toprule
Case & A & $v_0(A)$ & $\tau_A^4$ \\
\midrule
L=8, \(\ell{=}3\) & 3 & 0.5 & 7.983 \\
L=12, \(\ell{=}6\) & 4 & 0.2929 & — \\
L=12, \(\ell{=}6\) & 6 & 0.0349 & 8.708 \\
\bottomrule
\end{tabular}
\end{table}

\subsection{A6: Delta ladder TEBD}\label{app:a6}
No figure placeholder needed. Artifacts: \texttt{A6/A6\_summary.csv}, \texttt{A6/A6\_results.json}, \texttt{A6/esse\_delta\_ladder\_L12\_ell6\_k10.json}. Log: \texttt{A6/A6\_run.log}.

\paragraph{Summary (from JSON).}
L=12, \(\ell{=}6\), steps=150: commutator ladder fractions $=1.0$ at $\varepsilon\in\{10^{-8},10^{-6},10^{-4}\}$; norm-cap ladder $=0$ at $\alpha{=}1$; $\varepsilon_S\approx0.512$, $\varepsilon_I\approx1.025$, RHS$_{\mathrm{speed}}{=}0$.

\begin{table}[!ht]
\centering
\caption{A6: condensed metrics}
\label{tab:a6-condensed}
\begin{tabular}{@{}lccccc@{}}
\toprule
$\varepsilon_F$ & $\varepsilon_S$ & $\varepsilon_I$ & RHS$_{\mathrm{speed}}$ & RHS$_{\mathrm{cap}}$ & $\langle\widehat{\Psi}\rangle$ \\
\midrule
0.0 & 0.512 & 1.025 & 0.0 & 0.0798 & 0.0554 \\
\bottomrule
\end{tabular}
\end{table}

\subsection{A7: Commuting vs non-commuting comparison}\label{app:a7}


\paragraph{Summary (from log).}
Non-commuting: $\langle\Pi\rangle=1.546\pm0.285$, RHS$_{\mathrm{cap}}=0.107\pm0.021$, RHS$_{\mathrm{speed}}=0.0110\pm0.0020$, $\langle\widehat{\Psi}\rangle=0.0831\pm0.0163$.\\
Commuting: $\langle\Pi\rangle=0.841\pm0.087$, RHS$_{\mathrm{cap}}=0.0541\pm0.0055$, RHS$_{\mathrm{speed}}=0.00207\pm0.00209$, $\langle\widehat{\Psi}\rangle=0.0563\pm0.0058$.

\begin{table}[!ht]
\centering
\caption{A7: commuting vs non-commuting}
\label{tab:a7-compare}
\begin{tabular}{@{}lccccc@{}}
\toprule
Case & $\langle\Pi\rangle$ & $\tau_A^4$ & RHS$_{\mathrm{cap}}$ & RHS$_{\mathrm{speed}}$ & $\langle\widehat{\Psi}\rangle$ \\
\midrule
non-comm & $1.546\pm0.285$ & — & $0.107\pm0.021$ & $0.0110\pm0.0020$ & $0.0831\pm0.0163$ \\
comm     & $0.841\pm0.087$ & — & $0.0541\pm0.0055$ & $0.00207\pm0.00209$ & $0.0563\pm0.0058$ \\
\bottomrule
\end{tabular}
\end{table}

\subsection{A8: ETH eigenstate validation}\label{app:a8}


\paragraph{Summary (from log).}
N=8, A=[0,1]: LHS$=401.84$, RHS$_{\mathrm{cap}}=43.70$, RHS$_{\mathrm{speed}}=0$, $\langle\widehat{\Psi}\rangle=4.190$, $S_A=1.360$, $\sqrt{F_Q}=23.620$, $\Delta H_A=108.593$, $I=2.720$, $\Pi_A=87.403$.

\begin{table}[!ht]
\centering
\caption{A8: eigenstate metrics}
\label{tab:a8-eigen}
\begin{tabular}{@{}lccccc@{}}
\toprule
LHS & RHS$_{\mathrm{cap}}$ & RHS$_{\mathrm{speed}}$ & $\langle\widehat{\Psi}\rangle$ & $S_A$ \\
\midrule
401.84 & 43.70 & 0 & 4.190 & 1.360 \\
\bottomrule
\end{tabular}
\end{table}

\begin{table}[!ht]
\centering
\caption{A8: core metrics}
\label{tab:a8-core}
\begin{tabular}{@{}lccccc@{}}
\toprule
 & $\langle\Pi\rangle$ & $\tau_A^4$ & RHS$_{\mathrm{cap}}$ & RHS$_{\mathrm{speed}}$ & $\langle\widehat{\Psi}\rangle$ \\
\midrule
snapshot & — & — & 43.70 & 0 & 4.190 \\
\bottomrule
\end{tabular}
\end{table}

\subsection{A9: Ball-\(K_A\) proxy}\label{app:a9}
\begin{figure}[H]
  \centering
  \includegraphics[width=0.65\linewidth]{esse_ball_ka_L10_ell4_chi128_dt0.05_steps300_plot.png}
  \caption{A9: Conformal-like weighted $K_A$ proxy.}
  \label{fig:A9}
\end{figure}
\FloatBarrier
\noindent CSV: \texttt{A9/A9\_summary.csv}. Full traces: \texttt{A9/esse\_ball\_ka\_L10\_ell4\_chi128\_dt0.05\_steps300.csv}, 


\paragraph{Summary (from log).}
L=10, \(\ell{=}4\): $\Pi_{\mathrm{std}}^{\mathrm{avg}}=32.30$, $\Pi_{\mathrm{ball}}^{\mathrm{avg}}=71.17$, $\langle\widehat{\Psi}\rangle=0.0923$, LHS$=35.59$, RHS$_{\mathrm{cap}}=17.83$, RHS$_{\mathrm{speed}}=2.90$.

\begin{table}[!ht]
\centering
\caption{A9: summary}
\label{tab:a9}
\begin{tabular}{@{}lccccc@{}}
\toprule
$\Pi_{\mathrm{std}}$ & $\Pi_{\mathrm{ball}}$ & $\langle\widehat{\Psi}\rangle$ & LHS & RHS$_{\mathrm{cap}}$ & RHS$_{\mathrm{speed}}$ \\
\midrule
32.30 & 71.17 & 0.0923 & 35.59 & 17.83 & 2.90 \\
\bottomrule
\end{tabular}
\end{table}

\begin{table}[!ht]
\centering
\caption{A9: core metrics}
\label{tab:a9-core}
\begin{tabular}{@{}lccccc@{}}
\toprule
 & $\langle\Pi\rangle$ & $\tau_A^4$ & RHS$_{\mathrm{cap}}$ & RHS$_{\mathrm{speed}}$ & $\langle\widehat{\Psi}\rangle$ \\
\midrule
L=10, $\ell{=}4$ & — & — & 17.83 & 2.90 & 0.0923 \\
\bottomrule
\end{tabular}
\end{table}

\subsection{A10: Triad intersection density}\label{app:a10}
\begin{figure}[H]
  \centering
  \includegraphics[width=0.65\linewidth]{A10_triad_density.png}
  \caption{A10: Empirical lower density of triad intersection.}
  \label{fig:A10}
\end{figure}
\FloatBarrier
\noindent CSV: \texttt{A10/A10\_density.csv}. Log: \texttt{A10/A10\_run.log}.


\paragraph{Summary (from log).}
Empirical lower density: D(F)$\approx0.947$, D(S)$\approx1.000$, D(I)$\approx1.000$, D(all)$\approx0.947$ (lag-tolerant $\approx0.987$). LHS$\approx5.20$, RHS$_{\mathrm{cap}}\approx2.12$, RHS$_{\mathrm{speed}}\approx1.34\times10^{-3}$.

\begin{table}[!ht]
\centering
\caption{A10: intersection densities and thresholds}
\label{tab:a10}
\begin{tabular}{@{}lcccc@{}}
\toprule
 & D(F) & D(S) & D(I) & D(all) \\
\midrule
values & 0.947 & 1.000 & 1.000 & 0.947 \\
\midrule
 & LHS & RHS$_{\mathrm{cap}}$ & RHS$_{\mathrm{speed}}$ & \\
values & 5.20 & 2.12 & $1.34\times10^{-3}$ & \\
\bottomrule
\end{tabular}
\end{table}

\begin{table}[!ht]
\centering
\caption{A10: core metrics}
\label{tab:a10-core}
\begin{tabular}{@{}lccccc@{}}
\toprule
 & $\langle\Pi\rangle$ & $\tau_A^4$ & RHS$_{\mathrm{cap}}$ & RHS$_{\mathrm{speed}}$ & $\langle\widehat{\Psi}\rangle$ \\
\midrule
density case & — & — & 2.12 & $1.34\times10^{-3}$ & — \\
\bottomrule
\end{tabular}
\end{table}

\subsection{A11: Mixed-state quench stress}\label{app:a11}

\begin{figure}[H]
  \centering
  \includegraphics[width=0.75\linewidth]{stress_Pi1_vs_T_N4_hx0.5.png}
  \caption{A11: $N{=}4$, $h_x{=}0.5$.}
  \label{fig:A11_N4_hx05}
\end{figure}
\FloatBarrier

\begin{figure}[H]
  \centering
  \includegraphics[width=0.75\linewidth]{stress_Pi1_vs_T_N4_hx1.0.png}
  \caption{A11: $N{=}4$, $h_x{=}1.0$.}
  \label{fig:A11_N4_hx10}
\end{figure}
\FloatBarrier

\begin{figure}[H]
  \centering
  \includegraphics[width=0.75\linewidth]{stress_Pi1_vs_T_N6_hx0.5.png}
  \caption{A11: $N{=}6$, $h_x{=}0.5$.}
  \label{fig:A11_N6_hx05}
\end{figure}
\FloatBarrier

\begin{figure}[H]
  \centering
  \includegraphics[width=0.75\linewidth]{stress_Pi1_vs_T_N6_hx1.0.png}
  \caption{A11: $N{=}6$, $h_x{=}1.0$.}
  \label{fig:A11_N6_hx10}
\end{figure}
\FloatBarrier

\begin{figure}[H]
  \centering
  \includegraphics[width=0.75\linewidth]{stress_Pi1_vs_T_N8_hx0.5.png}
  \caption{A11: $N{=}8$, $h_x{=}0.5$.}
  \label{fig:A11_N8_hx05}
\end{figure}
\FloatBarrier

\begin{figure}[H]
  \centering
  \includegraphics[width=0.75\linewidth]{stress_Pi1_vs_T_N8_hx1.0.png}
  \caption{A11: $N{=}8$, $h_x{=}1.0$.}
  \label{fig:A11_N8_hx10}
\end{figure}
\FloatBarrier

\noindent CSV: \texttt{A11/A11\_mixed\_quench\_summary.csv}. Log: \texttt{A11/A11\_run.log}.

\paragraph{Summary (from log).}
Representative: $(N{=}4, A{=}1, T{=}0.5, h_x{=}1.0)$ LHS$=0.1286$, RHS$_{\mathrm{cap}}=0.0637$, RHS$_{\mathrm{speed}}=0.00930$, $\langle\widehat{\Psi}\rangle=0.0469$. 
Larger $A$ increases margins (e.g., $N{=}4, A{=}2, T{=}0.5, h_x{=}1.0$: LHS$=0.381$, RHS$_{\mathrm{cap}}=0.235$, RHS$_{\mathrm{speed}}=0.057$).

\begin{table}[H]
\centering
\caption{A11: Representative configurations.}
\label{tab:a11}
\begin{tabular}{@{}lcccc@{}}
\toprule
Config & LHS & RHS$_{\mathrm{cap}}$ & RHS$_{\mathrm{speed}}$ & $\langle\widehat{\Psi}\rangle$ \\
\midrule
N4,A1,T0.5,hx1.0 & 0.1286 & 0.0637 & 0.00930 & 0.0469 \\
N4,A2,T0.5,hx1.0 & 0.3813 & 0.2355 & 0.0571 & 0.0198 \\
\bottomrule
\end{tabular}
\end{table}
\FloatBarrier

\begin{table}[H]
\centering
\caption{A11: Core metrics (selected configurations).}
\label{tab:a11-core}
\begin{tabular}{@{}lccccc@{}}
\toprule
Config & $\langle\Pi\rangle$ & $\tau_A^4$ & RHS$_{\mathrm{cap}}$ & RHS$_{\mathrm{speed}}$ & $\langle\widehat{\Psi}\rangle$ \\
\midrule
N4,A1,T0.5,hx1.0 & — & — & 0.0637 & 0.00930 & 0.0469 \\
N4,A2,T0.5,hx1.0 & — & — & 0.2355 & 0.0571 & 0.0198 \\
\bottomrule
\end{tabular}
\end{table}
\FloatBarrier

\begin{table}[H]
\centering
\caption{A11: Condensed head/tail ($A{=}1$, $T\in\{0.5,2.0\}$).}
\label{tab:a11-condensed}
\begin{tabular}{@{}lcccc@{}}
\toprule
Config & $\langle\Pi\rangle$ & LHS & RHS$_{\mathrm{cap}}$ & RHS$_{\mathrm{speed}}$ \\
\midrule
N4, T0.5, hx0.5 & 0.0074 & 0.00472 & 0.00370 & 0.000453 \\
N4, T2.0, hx0.5 & 0.0153 & 0.05716 & 0.00764 & 0.001657 \\
N6, T0.5, hx0.5 & 0.0068 & 0.00438 & 0.00342 & 0.000471 \\
N6, T2.0, hx0.5 & 0.0165 & 0.05925 & 0.00824 & 0.001505 \\
N8, T0.5, hx0.5 & 0.0058 & 0.00373 & 0.00292 & 0.000470 \\
N8, T2.0, hx0.5 & 0.0168 & 0.05922 & 0.00840 & 0.001561 \\
\bottomrule
\end{tabular}
\end{table}
\FloatBarrier

\subsection{A12: Weak interaction sweep}\label{app:a12}
\begin{figure}[H]
  \centering
  \includegraphics[width=0.65\linewidth]{A12_weak_J_sweep.png}
  \caption{A12: Sweep over weak interaction strengths.}
  \label{fig:A12}
\end{figure}
\FloatBarrier
\noindent CSV: \texttt{A12/A12\_weak\_J\_sweep.csv}. Log: \texttt{A12/A12\_run.log}.

\paragraph{Summary (from log).}
N=8, J sweep $10^{-3}\to 1$: at J=0.001, LHS$\approx0.0284$, RHS$_{\mathrm{cap}}\approx0.0115$, RHS$_{\mathrm{speed}}\approx6.0\times10^{-4}$; at J=1.0, LHS$\approx0.775$, RHS$_{\mathrm{cap}}\approx0.359$, RHS$_{\mathrm{speed}}\approx0.175$, $\langle\widehat{\Psi}\rangle\approx0.0344$.

\begin{table}[!ht]
\centering
\caption{A12: endpoints of J-sweep}
\label{tab:a12}
\begin{tabular}{@{}lcccc@{}}
\toprule
J & LHS & RHS$_{\mathrm{cap}}$ & RHS$_{\mathrm{speed}}$ & $\langle\widehat{\Psi}\rangle$ \\
\midrule
0.001 & 0.0284 & 0.0115 & $6.0\times10^{-4}$ & 0.00139 \\
1.000 & 0.7750 & 0.3587 & 0.1752 & 0.03440 \\
\bottomrule
\end{tabular}
\end{table}

\begin{table}[!ht]
\centering
\caption{A12: condensed head/tail (J-sweep)}
\label{tab:a12-condensed}
\begin{tabular}{@{}lcccccc@{}}
\toprule
J & $\langle\Pi\rangle$ & LHS & RHS$_{\mathrm{cap}}$ & RHS$_{\mathrm{speed}}$ & $\langle\widehat{\Psi}\rangle$ \\
\midrule
0.001 & 0.0230 & 0.0284 & 0.0115 & $6.0\times10^{-4}$ & 0.00139 \\
1.000 & 0.7174 & 0.7749 & 0.3587 & 0.1752 & 0.03440 \\
\bottomrule
\end{tabular}
\end{table}

\subsection{A13: Integrable XX validation}\label{app:a13}


\begin{figure}[H]
  \centering
  \includegraphics[width=0.65\linewidth]{A13_xx_integrable.png}
  \caption{A13: XX-chain integrable case diagnostics and thresholds.}
  \label{fig:A13}
\end{figure}
\FloatBarrier
\noindent CSVs: \texttt{A13/A13\_xx\_integrable.csv}, \texttt{A13/A13\_thresholds.csv}. 
The thresholds table includes columns: \emph{N}, \emph{A\_size}, \emph{LHS\_energy\_triad}, 
\emph{RHS\_capped}, \emph{RHS\_speed}, \emph{psi\_hat\_mean}, \emph{delta\_speed}, 
\emph{H\_norm\_A}, \emph{dmin}. Log: \texttt{A13/A13\_run.log}.



\paragraph{Summary (from log).}
N=6, A=2: $\langle\Pi\rangle_T=7.351$, $\tau_A^4=6.286$, $v_0\approx0.559$, perC$_A\approx3.514$. Thresholds: LHS=3.675, RHS$_{\mathrm{cap}}=2.217$, RHS$_{\mathrm{speed}}=0.349$, $\langle\widehat{\Psi}\rangle=0.288$.

\begin{table}[!ht]
\centering
\caption{A13: core metrics}
\label{tab:a13-core}
\begin{tabular}{@{}lccccc@{}}
\toprule
 & $\langle\Pi\rangle$ & $\tau_A^4$ & RHS$_{\mathrm{cap}}$ & RHS$_{\mathrm{speed}}$ & $\langle\widehat{\Psi}\rangle$ \\
\midrule
N=6, A=2 & 7.351 & 6.286 & 2.217 & 0.349 & 0.288 \\
\bottomrule
\end{tabular}
\end{table}
\subsection{A14: Topological triad }\label{app:a14}
\begin{figure}[H]
  \centering
  \includegraphics[width=0.65\linewidth]{toric_topo_triad_L2_mixed-gs.png}
  \caption{A14: Toric code generalized triad for a torus ($L{=}2$), using the mixed ground space to expose logical-loop fluctuations.}
  \label{fig:A14_toric}
\end{figure}
\FloatBarrier
\noindent CSV: \texttt{A14/toric\_topo\_triad\_summary.csv}. JSON: \texttt{A14/toric\_topo\_triad\_L2\_mixed-gs.json}. Log: \texttt{A14/A14\_run.log}.

\begin{table}[!ht]
\centering
\caption{A14: L=2 (mixed ground space) core metrics}
\label{tab:a14}
\begin{tabular}{@{}lccccccc@{}}
\toprule
N & degeneracy & $S_A$ & $I(A{:}\bar A)$ & $\Delta_x$ & $\Delta_z$ & $\Pi_x$ & $\Pi_z$ \\
\midrule
8 & 4 & 2.773 & 4.159 & 1.000 & 1.000 & 11.531 & 11.531 \\
\bottomrule
\end{tabular}
\end{table}

\begin{table}[!ht]
\centering
\caption{A1: core metrics}
\label{tab:a1-core}
\begin{tabular}{@{}lccccc@{}}
\toprule
Case & $\langle\Pi_{\mathrm{rate}}\rangle$ & $\tau_A^4$ & RHS$_{\mathrm{cap}}$ & RHS$_{\mathrm{speed}}$ & $\langle\widehat{\Psi}\rangle$ \\
\midrule
non-comm & 1.157 & — & 0.0787 & 0.0150 & 0.061 \\
comm     & 0.870 & — & 0.0570 & $7.4\times10^{-5}$ & 0.059 \\
\bottomrule
\end{tabular}
\end{table}
\subsection*{Aggregated Annex Statistics}
For quick reference, we report representative summary metrics drawn from the CSV artifacts:
\begin{table}[!ht]
\centering
\caption{Aggregated Annex Statistics for selected runs (A1, A2, A5, A13).}
\label{tab:annex-agg}
\begin{tabular}{@{}lcccccc@{}}
\toprule
Label & Scenario & $\langle\Pi\rangle$ or $\langle\Pi_{\mathrm{rate}}\rangle$ & LHS (energy triad) & RHS$_{\mathrm{cap}}$ & RHS$_{\mathrm{speed}}$ & Notes \\
\midrule
A1 & non-comm & $1.157$ & $1.258$ & $0.079$ & $0.015$ & $N{=}6$ \\
A1 & comm & $0.870$ & $2.504$ & $0.057$ & $7.4\times 10^{-5}$ & $N{=}6$ \\
A2 & non-comm (trial 1) & $1.157$ & $1.258$ & $0.079$ & $0.015$ & $r\approx0.385$ \\
A2 & comm (trial 1) & $0.870$ & $2.504$ & $0.057$ & $\approx 0$ & $r\approx0.933$ \\
A5 & per-cut (TEBD) & $7.554$ & — & — & $1.284$ & $\varepsilon_F{=}1.303$ \\
A13 & XX integrable & — & $3.675$ & $2.217$ & $0.349$ & $\psi\_\mathrm{hat}{=}0.288$ \\
\bottomrule
\end{tabular}
\end{table}
All reported left-hand sides exceed both canonical right-hand-side thresholds in their contexts; intensive $\widehat{\Psi}$ means remain in $[0,1]$, consistent with normalization.

\paragraph{Extended summary (A1--A13).}
Key additional artifacts summarized from CSVs:
\begin{table}[!ht]
\centering
\caption{A3--A13: extended summary (standardized metrics; see per-run CSV/JSON).}
\label{tab:annex-extended}
\begin{tabular}{@{}lcccccc@{}}
\toprule
Label & Scenario & $\langle\Pi\rangle$ or $\langle\Pi_{\mathrm{rate}}\rangle$ & $\tau_A^4$ & RHS$_{\mathrm{cap}}$ & RHS$_{\mathrm{speed}}$ & $\langle\widehat{\Psi}\rangle$ \\
\midrule
A1 & non-comm & 1.157 & — & 0.0787 & 0.0150 & 0.061 \\
A1 & comm     & 0.870 & — & 0.0570 & $7.4\times10^{-5}$ & 0.059 \\
A2 & non-comm (trial 1) & 1.157 & — & 0.0787 & 0.0150 & — \\
A2 & comm (trial 1)    & 0.870 & — & 0.0570 & $7.4\times10^{-5}$ & — \\
A3 & 1-site (agg.) & — & — & $1.074\!\times\!10^{-1}$ & $1.162\!\times\!10^{-2}$ & 0.0831 \\
A4 & $\ell{=}3$, N=8 & 4.988 & 3.604 & 2.494 & 0.100 & 0.0560 \\
A4 & $\ell{=}4$, N=8 & 18.060 & 9.027 & 9.030 & 4.841 & 0.0868 \\
A5 & Run1 & 7.554 & — & 3.777 & 1.284 & — \\
A6 & L=12, $\ell{=}6$ & — & — & 0.0798 & 0.0 & 0.0554 \\
A7 & non-comm & $1.546\pm0.285$ & — & $0.107\pm0.021$ & $0.0110\pm0.0020$ & $0.0831\pm0.0163$ \\
A8 & snapshot & — & — & 43.70 & 0 & 4.190 \\
A9 & L=10, $\ell{=}4$ & — & — & 17.83 & 2.90 & 0.0923 \\
A10 & density & — & — & 2.12 & $1.34\times10^{-3}$ & — \\
A11 & N4,A1,T0.5,$h_x$1.0 & 0.1273 & 0.2888 & 0.0637 & 0.00930 & 0.0469 \\
A12 & J=1.0 & 0.7174 & — & 0.3587 & 0.1752 & 0.03440 \\
A13 & N=6, A=2 & 7.351 & 6.286 & 2.217 & 0.349 & 0.288 \\
\bottomrule
\end{tabular}
\end{table}

Where ranges indicate variation across the sweep or grid. Full details are available in the per-A CSVs listed above.

\begin{thebibliography}{99}

\bibitem{BraunsteinCaves1994}
S.~L. Braunstein and C.~M. Caves, 
“Statistical distance and the geometry of quantum states,” 
\emph{Phys. Rev. Lett.}, vol.~72, pp.~3439–3443, 1994.

\bibitem{NielsenChuang2010}
M.~A. Nielsen and I.~L. Chuang, 
\emph{Quantum Computation and Quantum Information}, 
Cambridge University Press, 2010.

\bibitem{Taddei2013}
M.~M. Taddei, B.~M. Escher, L. Davidovich, and R.~L. de Matos Filho, 
“Quantum speed limit for physical processes,” 
\emph{Phys. Rev. Lett.}, vol.~110, p.~050402, 2013.

\bibitem{DeffnerLutz2013}
S. Deffner and E. Lutz, 
“Quantum speed limit for non-Markovian dynamics,” 
\emph{Phys. Rev. Lett.}, vol.~111, p.~010402, 2013.

\bibitem{Blanco2013}
D. Blanco, H. Casini, L. Hung, and R. Myers, 
“Relative entropy and holography,” 
\emph{J. High Energy Phys.}, vol.~2013, no.~8, p.~060, 2013.

\bibitem{NachtergaeleSims2010}
B. Nachtergaele and R. Sims, 
“Lieb–Robinson bounds in quantum many-body physics,” 
\emph{Contemp. Math.}, vol.~552, pp.~161–175, 2011.

\bibitem{CsiszarKorner2011}
I. Csiszár and J. Körner, 
\emph{Information Theory: Coding Theorems for Discrete Memoryless Systems}, 
Cambridge University Press, 2011.

\bibitem{Safranek2017}
D. Šafránek, 
“Discontinuities of the quantum Fisher information and the Bures metric,” 
\emph{Phys. Rev. A}, vol.~95, p.~052320, 2017.

\bibitem{AlbaCalabrese2017}
V. Alba and P. Calabrese, 
“Entanglement and entropy production in quantum quenches: a quasiparticle picture,” 
\emph{Phys. Rev. B}, vol.~96, p.~115421, 2017.

\bibitem{Frowis2018}
F. Fröwis, P. Sekatski, W. Dür, N. Gisin, and N. Brunner, 
“Macroscopic quantum states: measures, witnesses, and uses,” 
\emph{Rev. Mod. Phys.}, vol.~90, p.~025004, 2018.

\bibitem{Spohn1978}
H. Spohn, 
“Entropy production for quantum dynamical semigroups,” 
\emph{J. Math. Phys.}, vol.~19, pp.~1227–1230, 1978.

\bibitem{BreuerPetruccione2002}
H.-P. Breuer and F. Petruccione, 
\emph{The Theory of Open Quantum Systems}, 
Oxford University Press, 2002.

\bibitem{Gorini1976}
V. Gorini, A. Kossakowski, and E.~C.~G. Sudarshan, 
“Completely positive dynamical semigroups of N-level systems,” 
\emph{J. Math. Phys.}, vol.~17, pp.~821–825, 1976.

\bibitem{Lindblad1976}
G. Lindblad, 
“On the generators of quantum dynamical semigroups,” 
\emph{Commun. Math. Phys.}, vol.~48, pp.~119–130, 1976.

\bibitem{Stinespring1955}
W.~F. Stinespring, 
“Positive functions on C$^*$-algebras,” 
\emph{Proc. Amer. Math. Soc.}, vol.~6, pp.~211–216, 1955.

\bibitem{Kliesch2014}
M. Kliesch, C. Gogolin, M.~J. Kastoryano, A. Riera, and J. Eisert, 
“Locality of temperature,” 
\emph{Phys. Rev. X}, vol.~4, p.~031019, 2014.

\bibitem{MandelstamTamm1945}
L. Mandelstam and I. Tamm, 
“The uncertainty relation between energy and time in nonrelativistic quantum mechanics,” 
\emph{J. Phys. USSR}, vol.~9, pp.~249–254, 1945.

\bibitem{MargolusLevitin1998}
N. Margolus and L.~B. Levitin, 
“The maximum speed of dynamical evolution,” 
\emph{Physica D}, vol.~120, pp.~188–195, 1998.

\bibitem{Horodecki2009}
R. Horodecki, P. Horodecki, M. Horodecki, and K. Horodecki, 
“Quantum entanglement,” 
\emph{Rev. Mod. Phys.}, vol.~81, p.~865, 2009.

\bibitem{BrandaoHorodecki2015}
F. Brandão and M. Horodecki, 
“Exponential decay of correlations implies area law,” 
\emph{Commun. Math. Phys.}, vol.~333, pp.~761–798, 2015.

\bibitem{Escher2011}
B. Escher, R. de Matos Filho, and L. Davidovich, 
“General framework for estimating the ultimate precision limit in noisy quantum-enhanced metrology,” 
\emph{Nat. Phys.}, vol.~7, pp.~406–411, 2011.

\bibitem{DeffnerCampbell2017}
S. Deffner and S. Campbell, 
“Quantum speed limits: from Heisenberg’s uncertainty principle to optimal quantum control,” 
\emph{J. Phys. A: Math. Theor.}, vol.~50, p.~453001, 2017.

\bibitem{Campisi2011}
M. Campisi, P. Hänggi, and P. Talkner, 
“Colloquium: Quantum fluctuation relations: Foundations and applications,” 
\emph{Rev. Mod. Phys.}, vol.~83, p.~771, 2011.

\bibitem{Hamazaki2024}
R. Hamazaki, 
“Speed limits to fluctuation dynamics,” 
\emph{Commun. Phys.}, vol.~7, p.~361, 2024. arXiv:2309.07301.

\bibitem{ShiraishiSaito2018}
N. Shiraishi and K. Saito, 
“Fundamental relation between entropy production and heat current,” 
\emph{J. Stat. Phys.}, vol.~174, pp.~433–453, 2019.

\bibitem{NishiyamaHasegawa2025}
T. Nishiyama and Y. Hasegawa, 
“Unified speed limits in classical and quantum dynamics via temporal Fisher information,” 
arXiv:2504.04790, 2025.

\bibitem{NishiyamaHasegawa2024}
T. Nishiyama and Y. Hasegawa, 
“Tradeoff relations in open quantum dynamics via Robertson, Maccone–Pati, and Robertson–Schrödinger uncertainty relations,” 
\emph{J. Phys. A: Math. Theor.}, vol.~57, p.~415301, 2024. arXiv:2402.09680.

\bibitem{BenattiFloreanini2005}
F. Benatti and R. Floreanini, 
“Open Quantum Dynamics: Complete Positivity and Entanglement,” 
\emph{Int. J. Mod. Phys. B}, vol.~19, p.~3063, 2005. arXiv:quant-ph/0507271.

\bibitem{Benedetti2012}
C. Benedetti, F. Buscemi, P. Bordone, and M. G. A. Paris, 
“Effects of classical environmental noise on entanglement and quantum discord,” 
\emph{Int. J. Quantum Inform.}, vol.~10, p.~1241005, 2012. arXiv:1209.4201.

\bibitem{CasiniHuertaMyers2011}
H. Casini, M. Huerta, and R.~C. Myers,
“Towards a derivation of holographic entanglement entropy,”
\emph{J. High Energy Phys.}, vol.~2011, no.~5, p.~036, 2011.

\bibitem{Petz1996}
D. Petz,
“Monotone metrics on matrix spaces,”
\emph{Linear Algebra Appl.}, vol.~244, pp.~81–96, 1996.

\bibitem{Hayashi2006}
M. Hayashi,
\emph{Quantum Information: An Introduction},
Springer, 2006.

\bibitem{FaulknerLewkowyczMaldacena2013}
T. Faulkner, A. Lewkowycz, and J. Maldacena,
“Quantum corrections to holographic entanglement entropy,”
\emph{J. High Energy Phys.}, vol.~2013, no.~11, p.~074, 2013.

\bibitem{Jacobson1995}
T. Jacobson,
“Thermodynamics of spacetime: The Einstein equation of state,”
\emph{Phys. Rev. Lett.}, vol.~75, pp.~1260–1263, 1995.

\bibitem{BrownHenneaux1986}
J.~D. Brown and M. Henneaux,
“Central charges in the canonical realization of asymptotic symmetries: An example from three-dimensional gravity,”
\emph{Commun. Math. Phys.}, vol.~104, pp.~207–226, 1986.

\end{thebibliography}

\end{document}

\end{document}